\documentclass[11pt]{article}
\usepackage[a4paper,margin=2.5cm]{geometry}
\usepackage{amsmath,amssymb,amsthm,mathtools}
\usepackage{enumitem}
\usepackage[colorlinks=true,linkcolor=blue,citecolor=blue,urlcolor=blue]{hyperref}

% ---------------------------------------------------------------- theorems
\newtheorem{theorem}{Theorem}[section]
\newtheorem{lemma}[theorem]{Lemma}
\newtheorem{proposition}[theorem]{Proposition}
\newtheorem{corollary}[theorem]{Corollary}
\newtheorem{claim}[theorem]{Claim}
\theoremstyle{definition}
\newtheorem{definition}[theorem]{Definition}
\newtheorem{convention}[theorem]{Convention}
\theoremstyle{remark}
\newtheorem{remark}[theorem]{Remark}

% ---------------------------------------------------------------- macros
\newcommand{\R}{\mathbb{R}}
\newcommand{\Sph}{\mathrm{S}}          % sphere S(c,t); S(c,U) union over radii U
\newcommand{\Bis}{\mathrm{H}}          % perpendicular bisector H(p,q)
\newcommand{\Dist}{\mathcal{D}}        % set of pairwise distances
\newcommand{\Bg}{\mathrm{Big}}
\newcommand{\rk}{\operatorname{rank}}
\newcommand{\aff}{\operatorname{aff}}
\newcommand{\dirsp}{\operatorname{dir}}
\newcommand{\Hyp}{\mathrm{Hyp}}
\newcommand{\Sm}{\mathrm{Sm}}
\newcommand{\Good}{\mathrm{Good}}
\newcommand{\Cov}{\mathrm{Cov}}
\newcommand{\Ex}{\mathrm{Ex}}
\newcommand{\Ty}{\mathrm{Ty}}
\newcommand{\Lst}{L^{\ast}}            % L_{k*}(D+2)
\newcommand{\kst}{k^{\ast}}
\newcommand{\smin}{\sigma_{\min}}
\newcommand{\dd}[2]{|#1#2|}            % distance |xy|
\newcommand{\binomc}[2]{\tbinom{#1}{#2}}
\newcommand{\eps}{\varepsilon}

\title{Subexponential bounds for distinct-distance subsets:\\
$f_d(n)=2^{O_n(\sqrt d\,\log d)}$ for every fixed $n$\\[4pt]
\large (an affirmative answer to Erd\H{o}s Problem \#1088, subject to expert verification)}

\author{%
\textbf{[PLACEHOLDER AUTHOR LIST --- to be confirmed by the team]}\\[4pt]
Will Howe (H0W4RD-dev), directing the work\\[2pt]
\small The mathematics was developed by AI agents (Claude, Anthropic) in sessions directed by the team;\\
\small see the statement on AI assistance and verification in Section~\ref{sec:ai}.}

\date{Draft manuscript, \today}

\begin{document}
\maketitle

\begin{abstract}
For integers $n\ge 3$ and $d\ge 1$ let $f_d(n)$ be the least $m$ such that every set of $m$ points in $\R^d$
contains $n$ points whose $\binom n2$ pairwise distances are all distinct. Erd\H{o}s asked whether
$f_d(n)=2^{o(d)}$ for every fixed $n$; the classical bound is exponential in $d$. We prove that
$f_d(n)\le 2^{C_n\sqrt d\,\log(d+2)}$ for every fixed $n\ge 3$, which would answer the question
affirmatively. The main new tool is a rank--rectangle theorem (Theorem~LR): if every entry of a real matrix of
rank $r$ equals one of $k$ prescribed functions, each depending only on the row or only on the column, then
there is a rectangle of relative size $2^{-O(k^3\sqrt{r}\log r)}$ in each direction on which a single one of these
functions describes all entries. Its proof combines Rothvoss's hyperplane-rounding proof of Lovett's
bound with a rescaling at the largest gap and a cleaning step. The geometric part of the argument is a
downward induction over ``levels'' (rainbow sets $S$ together with a ground set $W$ of points extending $S$),
with a polynomial-method kill step at each level and a structural class-count statement proved by repeated use
of Theorem~LR. \emph{Verification status:} every step was written out and checked by independent AI referee
passes. The whole argument has also been formalized in Lean~4 with Mathlib. The formalization proves the statement
\texttt{Erdos1088.erdos\_1088} of the Google DeepMind \emph{formal-conjectures} repository verbatim, using only
Lean's standard axioms (Section~\ref{sec:lean}). No human expert has reviewed the informal argument yet.
See Section~\ref{sec:ai}.
\end{abstract}

\tableofcontents

% ====================================================================
\section{Introduction}\label{sec:intro}
% ====================================================================

\subsection{The problem}
A finite set $A\subset\R^d$ is \emph{rainbow} if its $\binom{|A|}{2}$ pairwise distances are pairwise distinct
(sets with at most two points are rainbow). A set $X$ is \emph{rainbow-$n$-free} if it contains no rainbow
subset of size $n$. For $n\ge 3$ and $d\ge 1$ let $f_d(n)$ be the least $m$ such that every set of $m$ points of
$\R^d$ contains a rainbow $n$-subset, and let
\[
  s_n(j):=\sup\{|X| : X\subset\R^{N}\text{ finite, rainbow-$n$-free, } \dim\aff X\le j\}
\]
(the ambient dimension $N$ is irrelevant). Then $f_d(n)=s_n(d)+1$, and $s_n(0)=1$.

Erd\H{o}s asked whether $f_d(n)=2^{o(d)}$ for every fixed $n$; this is Problem~\#1088 on the
Erd\H{o}s problems website~\cite{EP1088}.\footnote{Within the project the problem is also tracked under
the identifier JSP-000905.} We prove that the answer is yes, with a much stronger quantitative bound.

\begin{theorem}[Main Theorem]\label{thm:main-intro}
For every fixed $n\ge 3$ there is a constant $C_n$ such that every set of more than
$2^{C_n\sqrt d\,\log_2(d+2)}$ points of $\R^d$ contains $n$ points with all $\binom n2$ pairwise distances
distinct. That is, $f_d(n)=2^{O_n(\sqrt d\log d)}=2^{o(d)}$.
\end{theorem}

\subsection{History}
As recorded on the problem page~\cite{EP1088}, Erd\H{o}s and Straus showed $f_d(n)\le c_n^{\,d}$ for a constant
$c_n$; the problem asks whether the exponential dependence on $d$ is necessary. In the other direction,
$s$-distance sets with $s=\binom n2-1$ are rainbow-$n$-free, since a rainbow $n$-set needs $\binom n2$ distinct
distances. For instance the $0/1$-vectors of weight $w$ in $\R^{d+1}$ lie in a $d$-dimensional affine subspace
and determine only the distances $\sqrt{2i}$, $1\le i\le w$; with $w=\binom n2-1$ this gives
$f_d(n)>\binom{d+1}{\binom n2-1}$ for $d+1\ge \binom n2 -1$. So $f_d(n)$ grows at least polynomially in $d$, with
degree depending on $n$, and we are not aware of any super-polynomial lower bound.

For $n=3$ the problem is classical and polynomial: a set is rainbow-$3$-free exactly when every three of its
points span an isosceles triangle, where collinear triples with two equal distances count as (degenerate)
isosceles triangles. Blokhuis~\cite{Blokhuis} proved that an isosceles set in $\R^d$ has at most
$\binom{d+2}{2}$ points. His argument is purely metric, so it covers degenerate triples, and it gives
$f_d(3)\le\binom{d+2}2+1$. For $n\ge 4$ we know of no bound of the form $2^{o(d)}$ in the
literature before the present work.

\subsection{The new tool: Theorem LR}
The geometric part of our argument repeatedly produces matrices of the form $M[x,w]=|x-w|^2$, of rank at most
$D+2$, in which every entry is ``explained'' by one of boundedly many simple causes. A typical cause is
$|x-w|=|w-t|$ for a fixed point $t$, which depends only on the column $w$. Another is $|x-w|=|x-t|$, which
depends only on the row $x$. A third is $|x-w|=d$ for a fixed number $d$. The tool that converts such a cover
into geometric structure is the following rank--rectangle theorem (Theorem~\ref{thm:LR}).

\medskip\noindent\textbf{Theorem LR (informal).} \emph{Let $M\in\R^{m\times n}$ have rank at most $r$, and let
$\tau_1,\dots,\tau_k$ be functions on $[m]\times[n]$, each depending only on the row or only on the column. If every
entry satisfies $M[x,w]=\tau_i(x,w)$ for some $i$, then there are $R\subseteq[m]$, $C\subseteq[n]$ and a single
$i$ with $M=\tau_i$ on $R\times C$, $|R|\ge 2^{-L}m$ and $|C|\ge 2^{-L}n$, where
$L=O(k^3\sqrt{r+2}\,\log(r+2))$.}

\medskip
For Boolean matrices this is the rectangle form of Lovett's bound~\cite{Lovett}, in Rothvoss's
formulation~\cite{Rothvoss}. Our proof follows Rothvoss's hyperplane-rounding argument and adds two ingredients.
The first is a normal form, obtained on a sub-rectangle of constant relative size, which allows us to rescale the
matrix at its largest gap \emph{by a row function, a column function or a constant}. This keeps a margin even
when the type values accumulate. The second is a cleaning step, the Gavinsky--Lovett lemma (Lemma~\ref{lem:GL}),
applied to a Hadamard product, whose rank enters only logarithmically.

\begin{remark}[possible independent interest; to be confirmed by experts]\label{rem:SS}
The special case in which all types are column functions concerns \emph{$k$-listable} matrices, in which every
column takes at most $k$ distinct values. Here Theorem~LR yields a column-constant submatrix of relative size
$2^{-O(k^3\sqrt d\log d)}$ in each direction for rank-$d$ matrices (Corollary~\ref{cor:RQlist}). By
Singer--Sudan~\cite[Prop.~2.3]{SingerSudan}, a parallel $k$-partitioned point--hyperplane configuration in $\R^d$
yields a $k$-listable matrix of rank at most $d$. A column-constant (``$1$-listable'') submatrix of it then yields a
complete bipartite subgraph of the incidence graph, of proportional size up to a factor depending on $k$. Singer and
Sudan write that ``the best bound we know for parallel 3-partitioned configurations is $2^{-\tilde O(d)}$''
(\cite{SingerSudan}, arXiv version~5, Section~1). As far as we can tell, Corollary~\ref{cor:RQlist} improves this to
$2^{-O(\sqrt d\log d)}$ for every fixed $k$. Our literature search (summarised in Section~\ref{sec:concl}) was
limited, and the observation follows from Rothvoss's proof in a few lines, so experts may regard it as known. The
claim of improvement must be confirmed by experts before it is relied on. Theorem~LR does not bear on the
log-rank conjecture, which would require $2^{-\mathrm{polylog}(d)}$.
\end{remark}

\subsection{Prior in-project work}
This manuscript builds on earlier notes of a peer session in the same project~\cite{Peer}. Those notes
recorded the problem, its literature and the lower-bound constructions above. They also contained
polynomial-method ``peeling'' lemmas: a point all of whose distances to the rest of a set lie in a short list can
be removed using a low-degree polynomial in the lifted coordinates $(|y|^2,y)$. Lemma~H below
(Lemma~\ref{lem:H}) is a two-sided variant of those peeling lemmas, in which each pair needs to be ``killed'' from
one side only, at the price of squaring the dimension. The present argument does not use the other results of
those notes, which include reductions to prime sets and twin analyses.

\subsection{Verification status}
We state this plainly, because the result would resolve a problem of Erd\H{o}s and the argument is long.
\begin{itemize}[leftmargin=2em]
\item The mathematics was developed by AI agents (Claude, Anthropic) in sessions directed by the team.
\item Each component was checked by at least one independent AI referee pass. These passes found and repaired
  one genuine gap, a threshold gap in the top-level argument (fix F1, incorporated in
  Section~\ref{sec:top}), together with a number of expository gaps, all addressed here.
\item \textbf{Machine-checked.} The complete argument has been formalized in Lean~4 (Section~\ref{sec:lean}).
  The formal proof establishes the statement of Problem~\#1088 in the formal-conjectures repository verbatim. It
  depends only on the axioms \texttt{propext}, \texttt{Classical.choice} and \texttt{Quot.sound}, and it has no
  \texttt{sorry}. The formalization proves Theorem~LR with a weaker dependence on~$k$ than stated here; this suffices
  for the Main Theorem.
\item \textbf{No human expert has reviewed the informal argument.} Theorem~LR is new and unpublished, and every later
  step depends on it. It is the component we most want checked by specialists.
\item This manuscript consolidates the project notes. A few routine changes made in the consolidation, listed in
  Section~\ref{sec:ai}, have not themselves been refereed; among them are some explicit constants and the direct
  single-cap write-up of Theorem~R$^+$.
\item The finite, threshold-free implications were also tested by computer in exact arithmetic. Because every
  threshold in the argument is astronomically large, such tests are sanity checks only and cannot test the
  asymptotic claims (Section~\ref{sec:concl}).
\end{itemize}
Until an expert check has been carried out, the Main Theorem should be regarded as a \emph{candidate} solution.

\subsection{Organisation}
Section~\ref{sec:structure} gives an overview of the proof, the linear order of dependencies and the threshold
tower. Section~\ref{sec:prelim} collects the classical tools: John's theorem, Sheppard's formula, the
cleaning lemma (with proof), Lemma~H, the Caro--Wei and K\"onig theorems, and the orthogonal split. Section~\ref{sec:LR}
proves Theorem~LR. Section~\ref{sec:G1} proves the label lemma and its consequence~G1 for pairs with disjoint label
sets. Section~\ref{sec:framework} sets up the level-configurations, the hereditary relative bound~HR and the kill
step. Section~\ref{sec:top} treats the top three levels. Sections~\ref{sec:core} and~\ref{sec:slicing} prove the
class-count statement at the lower levels: the Core Lemma and the Slicing Lemma. Section~\ref{sec:main} closes the
induction and proves the Main Theorem. Section~\ref{sec:concl} contains concluding remarks and the computational
sanity checks, and Section~\ref{sec:ai} the statement on AI assistance and the acknowledgements.

% ====================================================================
\section{Structure of the proof}\label{sec:structure}
% ====================================================================

Throughout, $n\ge 4$ is fixed; the case $n=3$ is Blokhuis's theorem. $X$ is a finite rainbow-$n$-free set of
affine dimension at most $D\ge 1$, and $h:=s_n(\lfloor D/2\rfloor)$. The proof is a \emph{halving recursion}: we
show
\begin{equation}\label{eq:halving-intro}
  s_n(D)\;\le\;2^{C_n\sqrt D\log_2(D+2)}\,\bigl(s_n(\lfloor D/2\rfloor)+1\bigr)\qquad(D\ge 1),
\end{equation}
and iterating~\eqref{eq:halving-intro} gives $s_n(d)=2^{O_n(\sqrt d\log d)}$ (Lemma~\ref{lem:halving}).
The quantity $h$ enters only through the \emph{orthogonal split} (Lemma~\ref{lem:OR}). If $R,C\subseteq X$ and
$\langle x-x',w-w'\rangle=0$ for all $x,x'\in R$ and $w,w'\in C$, then one of $R,C$ has affine dimension at most
$\lfloor D/2\rfloor$, so $\min(|R|,|C|)\le h$. Every rectangle produced by Theorem~LR in the geometric part ends
in such an orthogonal split.

\subsection{Levels}
A \emph{level-$j$ configuration} (Definition~\ref{def:config}) is a triple $(S,W,\sigma)$. Here $S\subseteq X$ is
a rainbow set with $|S|=j$, and $W\subseteq X\setminus S$ consists of points $x$ with $S\cup\{x\}$ rainbow. The
number $\sigma$ is a \emph{cap}: every sphere centred at a point of $S$ meets $W$ in fewer than $\sigma$
points. The statement
\begin{center}
\textbf{Claim$(j)$:} every level-$j$ configuration with $\sigma\ge\smin$ has $|W|\le\Gamma_j\,\sigma$
\end{center}
is proved by downward induction on $j$, from $j=n-2$ to $j=0$. Claim$(0)$, applied to $(\emptyset,X,\smin)$,
gives $|X|\le\Gamma_0\smin$. Here $\Gamma_0=(D+2)^{O_n(1)}$ and $\smin=2^{O_n(\sqrt D\log D)}(h+1)$, which
is~\eqref{eq:halving-intro}.

\subsection{The induction step}
Each step from Claim$(j+1)$ to Claim$(j)$ follows one pattern.
\begin{enumerate}[label=(\roman*),leftmargin=2.2em]
\item \emph{Kill step} (Lemma~\ref{lem:kill}, Theorem~\ref{thm:Rplus}). Each $b\in W$ receives a
  polynomial $F_b$ of bounded degree in the lifted coordinates $(|y|^2,y)$. Its factors are spheres about $b$
  with the radii of the \emph{big} classes of $b$, together with some degenerate loci. A point $x$ with
  $F_b(x)\ne 0$ extends $S\cup\{b\}$ to a rainbow set and lies only in small classes of $b$, so these points
  form a level-$(j+1)$ configuration and are few by Claim$(j+1)$. Caro--Wei and Lemma~H (a pair needs to be
  killed from one side only) then bound $|W|$. This works as long as every $b\in W$ has at most $K=O_n(1)$ big
  classes.
\item \emph{Few big classes.} At the top levels $j\ge n-4$ this is not needed. If $b$ has two big classes, then
  $n-4$ points of $S$ serve as ``spread'' points: together with $b$ and an auxiliary point $a$ they form an
  $(n-2)$-set. A label-counting argument (G1) and Theorem~LR then show that the auxiliary points $a$ which are not
  killed are few (Section~\ref{sec:top}; this is Theorem~R$^+$, written out in full for a single cap). At the
  lower levels $j=r\le n-5$ there are not enough spread points. Instead we prove the class-count statement
  \textbf{E1$_r$}: every $b\in W$ has fewer than $K_r=12\cdot 6^{\,n-5-r}$ big classes (Theorem~\ref{thm:E1}).
\item \emph{Proof of E1$_r$.} Suppose $b$ has $K_r$ big classes. We pass to a \emph{rainbow family}: centres
  $B\ni b$ and pairwise disjoint classes $C_k$ on which every point of $B$ has a constant distance, with all
  these constant labels distinct. The \emph{Slicing Lemma} (Section~\ref{sec:slicing}) adds one centre at a time
  and loses a factor $6$ in the number of classes. Its proof uses K\"onig's theorem, the relative bound HR and
  Theorem~LR. After $n-5-r$ slicing rounds there are $n-4-r$ centres and $12$ classes. The \emph{Core Lemma}
  (Section~\ref{sec:core}) shows that this is impossible, using two auxiliary points, G1, HR and three applications
  of Theorem~LR.
\item \emph{Hereditary relative bound HR$_r$} (Lemma~\ref{lem:HR}). It is a direct consequence of
  Claim$(r+1)$: for $c\in W$, $V\subseteq W\setminus\{c\}$ and $\theta\ge\sigma$, the points of $V$ lying in
  classes of $c$ that meet $V$ in fewer than $\theta$ points number at most $\Gamma'_r\theta$, with the explicit
  constant $\Gamma'_r=\Gamma_{r+1}+r^2$. This is the only place where the induction hypothesis enters the
  proof of E1$_r$, and it turns ``$c$ has few heavy classes'' into a \emph{covering} statement, which is what the
  cleaning lemma needs.
\end{enumerate}

\subsection{Linear order of dependencies}
The results are proved in the following order. Each item uses only earlier items, and none uses the
statement it is proving (there is no circularity).
\begin{enumerate}[label=\arabic*.,leftmargin=2em]
\item Classical tools (Section~\ref{sec:prelim}). John's theorem and Sheppard's formula (cited); the cleaning
  lemma GL (proved); the Hadamard-product rank bound HP and the dimension count DIM (proved); Lemma~H (proved);
  Caro--Wei and K\"onig (cited); the orthogonal split OR (proved); the halving lemma (proved).
\item Theorem~LR (Section~\ref{sec:LR}). Lemma~R (margin rounding), Lemmas~A and~B, Lemma~NF (normal form), then
  Theorem~LR, followed by Corollary~\ref{cor:RQlist}.
\item The label lemma and G1 (Section~\ref{sec:G1}), from LR and OR.
\item Configurations; Lemma~HR$_r$ and the kill Lemma~\ref{lem:kill}, both assuming Claim$(r+1)$
  (Section~\ref{sec:framework}).
\item Theorem~R$^+$: Claim$(j)$ for $j=n-2,n-3,\dots,\max(n-4,0)$ (Section~\ref{sec:top}), from G1, LR, GL, HP,
  OR, Lemma~H, Caro--Wei, DIM and the two-distribution Lemma~\ref{lem:twodist}.
\item For $n\ge 5$: the Core Lemma (Section~\ref{sec:core}) and the Slicing Lemma (Section~\ref{sec:slicing}),
  each assuming HR$_r$; then Theorem~E1$_r$.
\item The downward induction $\text{Claim}(r+1)\Rightarrow\text{HR}_r\Rightarrow\text{E1}_r\Rightarrow
  \text{Claim}(r)$ for $r=n-5,\dots,0$ (Section~\ref{sec:main}).
\item The Main Theorem, from Claim$(0)$ and the halving lemma (Section~\ref{sec:main}).
\end{enumerate}
The induction over levels takes place at a fixed dimension $D$, and the only input from dimensions below $D$ is
the finite number $h=s_n(\lfloor D/2\rfloor)$. The recursion~\eqref{eq:halving-intro} is therefore an ordinary
induction on $D$.

\subsection{The threshold tower}
All constants depend only on $n$ and $D$, except $\smin$, which also depends on $h$. They are defined in the
following order, and each definition uses only quantities defined before it (Section~\ref{sec:main} gives the
details).
\begin{enumerate}[label=(T\arabic*),leftmargin=2.6em]
\item Type counts $k'$, $k_c(r)$, $k_{\mathrm{sl}}(r,p)$; the class counts $K_r$; their maximum $\kst$. These
  depend on $n$ only.
\item $\Lst:=L_{\kst}(D+2)$ from Theorem~LR; then $\mu:=2^{\Lst}(h+1)$, $e:=2^{\Lst+1}(h+1)$,
  $e_A:=2^{\Lst+2}(h+1)$, and
  \[\smin:=2^{\Lst+2+2n}(h+1).\]
  So $\smin$ depends only on $n$, $D$ and $h$, and not on any $\Gamma$.
\item $Q:=2^{44}n^4(D+2)^{8k'}$ and $\Gamma_{n-2}:=4(n-2)^2$; then $\Gamma_{n-3}$ and $\Gamma_{n-4}$
  by~\eqref{eq:Gamma-top} (for $n=4$ these are $\Gamma_1,\Gamma_0$, and the tower ends here).
\item For $r=n-5,n-6,\dots,0$, in this order: $\Gamma'_r:=\Gamma_{r+1}+r^2$, then the ratios
  $\theta_{\mathrm{rel}}$, $\eps$, $\hat M_{\mathrm{core}}$, $\theta_{\mathrm{sl}}$, $\hat m_p$ and ${\mathrm{ratio}}_r$ (all
  functions of $n$, $D$ and $\Gamma'_r$), and then $\Gamma_r$ by~\eqref{eq:Gamma-low}.
\end{enumerate}
Every $\Gamma_j$, $\Gamma'_r$ and ${\mathrm{ratio}}_r$ is a pure number, at most $(D+2)^{c_n}$, and does not depend
on the cap $\sigma$. Every absolute error term occurring in the proofs ($\mu$, $e$, $e_A$, $2\cdot 4^{n-2}\mu+4$,
$4^n\mu$, and $1$) is at most $\smin\le\sigma$. These are the two facts that allow the caps to vary freely
inside the induction: Claim$(j+1)$ is applied with caps $\theta\ge\sigma$, $Q\sigma$ and
${\mathrm{ratio}}_r\sigma$, all of which are at least $\smin$.

% ====================================================================
\section{Notation, conventions and classical tools}\label{sec:prelim}
% ====================================================================

\subsection{Notation}
For $x,y\in\R^D$ we write $\dd xy:=|x-y|$. For $c\in\R^D$ and $t>0$, $\Sph(c,t):=\{y:\dd yc=t\}$ is the sphere
with centre $c$ and radius $t$, and for a set $U$ of positive reals $\Sph(c,U):=\bigcup_{u\in U}\Sph(c,u)$. For
$p\ne q$, $\Bis(p,q):=\{y:\dd yp=\dd yq\}$ is the perpendicular bisector hyperplane. For a finite set $A$,
$\Dist(A):=\{\dd xy: x\ne y\in A\}$ is its set of distances. If $A$ is rainbow then $|\Dist(A)|=\binom{|A|}2$. If
$c$ is a point and $V$ a set, a \emph{class of $c$ in $V$} is a nonempty set $V\cap\Sph(c,t)$ with $t>0$. Two
spheres about the same centre with different radii are disjoint, so the classes of $c$ in $V$ partition
$V\setminus\{c\}$.

The \emph{lifted coordinates} of $y\in\R^D$ are $(|y|^2,y_1,\dots,y_D)$. Every sphere equation
$\dd yc^2-t^2=|y|^2-2\langle y,c\rangle+|c|^2-t^2$ and every bisector equation
$\dd yp^2-\dd yq^2=-2\langle y,p-q\rangle+|p|^2-|q|^2$ is an affine function of the lifted coordinates.

Since an affine isometry changes neither distances nor affine dimension, we may and do assume that a set of
affine dimension at most $D$ lies in $\R^D$.

\begin{convention}[rectangles and rounding]\label{conv:rounding}
A \emph{rectangle} of a matrix with row set $A$ and column set $B$ is a product $A'\times B'$ with $A'\subseteq A$
and $B'\subseteq B$. Every rectangle produced below is nonempty, and its size bounds are real inequalities: a
conclusion $|A'|\ge 2^{-L}|A|$ with $2^{-L}|A|<1$ just says that $A'\ne\emptyset$. In the cleaning lemma,
``half'' means at least $|A|/2$ (hence at least $\lceil |A|/2\rceil$). Whenever a bound on a set $Z$ is deduced
from a bound on a half $Z''\subseteq Z$, we use $|Z|\le 2|Z''|$.
\end{convention}

\subsection{Two classical facts used in the rounding lemma}

\begin{theorem}[John~\cite{John}; see also~\cite{Ball}]\label{thm:John}
If $K\subset\R^s$ is a compact convex body that is symmetric about $0$ and has nonempty interior, then there is
an ellipsoid $E$ centred at $0$ with $E\subseteq K\subseteq\sqrt s\,E$.
\end{theorem}

\begin{theorem}[Sheppard's formula~\cite{Sheppard}; see also~\cite{GW}]\label{thm:Sheppard}
Let $u,v\in\R^s$ be unit vectors with $\langle u,v\rangle=a$ and let $g\sim N(0,I_s)$. Then
$\Pr[\langle g,u\rangle\ge0\text{ and }\langle g,v\rangle\ge0]=\tfrac12-\tfrac{\arccos a}{2\pi}
=\tfrac14+\tfrac{\arcsin a}{2\pi}$.
\end{theorem}
Consequently this probability is at least $\tfrac14$ if $a\ge 0$. If $a\le -b$ with $0\le b\le 1$, it is at
most $\tfrac14-\tfrac{\arcsin b}{2\pi}\le\tfrac14(1-\tfrac{2b}{\pi})$, because $\arcsin b\ge b$.

\subsection{Rank tools}

\begin{lemma}[Hadamard products, HP]\label{lem:HP}
For matrices $A,B$ of the same shape, $\rk(A\circ B)\le\rk A\cdot\rk B$, where $\circ$ is the entrywise
product.
\end{lemma}
\begin{proof}
Write $A=\sum_{i\le\rk A}a_ib_i^{T}$ and $B=\sum_{j\le\rk B}c_jd_j^{T}$. Then
$A\circ B=\sum_{i,j}(a_i\circ c_j)(b_i\circ d_j)^{T}$.
\end{proof}

\begin{lemma}[distance matrices]\label{lem:distrank}
Let $x_1,\dots,x_m,w_1,\dots,w_{m'}\in\R^D$ and let $f,g$ be real functions on the index sets. The matrix
$[\,\dd{x_i}{w_j}^2-f(i)-g(j)\,]_{i,j}$ has rank at most $D+2$.
\end{lemma}
\begin{proof}
It equals $(|x_i|^2-f(i))\mathbf 1^{T}+\mathbf 1(|w_j|^2-g(j))^{T}-2[\langle x_i,w_j\rangle]$, a sum of matrices of
rank at most $1$, $1$ and $D$.
\end{proof}
In particular, each of the following matrices has rank at most $D+2$: $[\dd xa^2-v^2]$ for a constant $v$;
$[\dd xa^2-v(a)^2]$ for a column function $v$; and $[\dd xa^2-\dd xs^2]$ for a fixed point $s$ (take
$f(x)=\dd xs^2$). By Lemma~\ref{lem:HP}, a Hadamard product of $k$ such matrices has rank at most $(D+2)^k$.

\begin{lemma}[cleaning lemma, GL; Gavinsky--Lovett~\cite{GavinskyLovett}, cf.~{\cite[Lemma 4.7]{HMST}}]
\label{lem:GL}
Let $B$ be a matrix over a field, with row set $A$ ($|A|=m$) and column set $B_0$ ($|B_0|=n$), and with
$\rk B\le R$ for some $R\ge 1$. Suppose $B$ has at most $mn/(16R)$ nonzero entries. Then there are
$A''\subseteq A$ and $B''\subseteq B_0$ with $|A''|\ge 7m/8$, $|B''|\ge n/2$ and $B=0$ on $A''\times B''$.
\end{lemma}
\begin{proof}
Call a column \emph{light} if it has at most $m/(8R)$ nonzero entries. If $c$ columns are heavy then
$c\,m/(8R)<mn/(16R)$ whenever $c\ge1$, so $c<n/2$, and the set $B''$ of light columns has $|B''|\ge n/2$.
The columns in $B''$ span a space of dimension at most $R$. Choose a basis of it consisting of at most $R$
columns from $B''$. The union of their supports has at most $R\cdot m/(8R)=m/8$ rows; let $A''$ be the
remaining rows. Every basis column vanishes on $A''$, and every column of $B''$ is a linear combination of basis
columns, so $B=0$ on $A''\times B''$.
\end{proof}

We use Lemma~\ref{lem:GL} in two forms. In the \emph{average form}, the total number of nonzero entries is at most
$mn/(16R)$. In the \emph{per-column form}, every column has at most $m/(16R)$ nonzero entries, which implies the
average form.

\begin{lemma}[dimension count, DIM]\label{lem:DIM}
The polynomials of degree at most $e$ in the $D+1$ lifted coordinates form a vector space of dimension
$\binom{D+1+e}{e}\le(D+2)^e$. A product of $e$ affine functions of the lifted coordinates lies in this space.
\end{lemma}
\begin{proof}
The count of monomials of degree at most $e$ in $D+1$ variables is $\binom{D+1+e}{e}=\prod_{i=1}^e\frac{D+1+i}{i}$,
and each factor is at most $D+2$.
\end{proof}

\begin{lemma}[Lemma H, symmetric form]\label{lem:H}
Let $I$ be a finite set and, for each $x\in I$, let $F_x$ be a real function on $I$. Suppose all $F_x$ lie in a
vector space of functions of dimension at most $m$, that $F_x(x)\ne0$ for all $x$, and that for every pair
$x\ne y$ in $I$ we have $F_x(y)=0$ or $F_y(x)=0$. Then $|I|\le m^2$.
\end{lemma}
\begin{proof}
Let $A_{y,x}:=F_y(x)$. Its rows are restrictions of functions from an $m$-dimensional space, so $\rk A\le m$. The
matrix $A\circ A^{T}$ has entries $F_y(x)F_x(y)$, which vanish for $x\ne y$ and are nonzero for $x=y$. It is
therefore a nonsingular diagonal matrix of size $|I|$. By Lemma~\ref{lem:HP}, $|I|=\rk(A\circ A^{T})\le m^2$.
\end{proof}
Lemma~H is the symmetric special case of a ``bi-peeling'' lemma from earlier in the project, which extends the
one-sided peeling lemmas of~\cite{Peer}. The squaring of the dimension is harmless for us.

\subsection{Graph theory}

\begin{theorem}[Caro--Wei~\cite{Caro,Wei}]\label{thm:CW}
A finite graph with $N$ vertices and average degree $\bar d$ has an independent set of size at least
$N/(\bar d+1)$.
\end{theorem}

\begin{theorem}[K\"onig~\cite{Konig}]\label{thm:Konig}
In a finite bipartite graph, the maximum size of a matching equals the minimum size of a vertex cover.
\end{theorem}

\subsection{The orthogonal split}

\begin{lemma}[orthogonal split, OR]\label{lem:OR}
Let $X\subset\R^D$ be rainbow-$n$-free, $h=s_n(\lfloor D/2\rfloor)$, and let $R,C\subseteq X$ be nonempty.
Suppose one of the following holds:
\begin{enumerate}[label=(\alph*)]
\item every $w\in C$ is equidistant from all points of $R$;
\item there is a point $p\notin C$ with $R\subseteq\Bis(p,w)$ for every $w\in C$.
\end{enumerate}
Then $\langle x-x',w-w'\rangle=0$ for all $x,x'\in R$ and $w,w'\in C$, and $\min(|R|,|C|)\le h$.
\end{lemma}
\begin{proof}
(a) $\dd xw^2-\dd{x'}w^2=|x|^2-|x'|^2-2\langle x-x',w\rangle$ does not depend on $w\in C$. Subtracting the
values at $w$ and $w'$ gives $\langle x-x',w-w'\rangle=0$. (b) $x\in\Bis(p,w)$ means
$2\langle x,w-p\rangle=|w|^2-|p|^2$. Hence $\langle x-x',w-p\rangle=0$ for $x,x'\in R$ and $w\in C$, and
subtracting the identities for $w$ and $w'$ gives the claim.
Now let $U,V\subseteq\R^D$ be the direction spaces of $\aff R$ and $\aff C$. They are orthogonal, so
$\dim U+\dim V\le D$ and $\min(\dim U,\dim V)\le\lfloor D/2\rfloor$. The corresponding set ($R$ or $C$) is a
rainbow-$n$-free set of affine dimension at most $\lfloor D/2\rfloor$, so its size is at most $h$.
\end{proof}

\subsection{The halving recursion}

\begin{lemma}[halving]\label{lem:halving}
Let $s:\mathbb Z_{\ge0}\to\mathbb Z_{\ge1}$ satisfy $s(0)=1$ and
$s(D)\le A(D)\,(s(\lfloor D/2\rfloor)+1)$ for all $D\ge1$, where $A(D)=2^{C\sqrt D\log_2(D+2)}$ with $C\ge1$.
Then $s(D)\le 2^{6C\sqrt D\log_2(D+2)}$ for all $D\ge1$.
\end{lemma}
\begin{proof}
Put $u(D):=\log_2(s(D)+1)$. Then $u(0)=1$, and for $D\ge 1$,
$u(D)\le\log_2\bigl(2A(D)(s(\lfloor D/2\rfloor)+1)\bigr)=1+C\sqrt D\log_2(D+2)+u(\lfloor D/2\rfloor)$.
Let $D_i:=\lfloor D/2^i\rfloor$ and $m:=\lfloor\log_2D\rfloor$, so that $D_m\ge1$ and $D_{m+1}=0$. Iterating, and
using $\sqrt{D_i}\le 2^{-i/2}\sqrt D$,
\[
u(D)\le (m+1)+1+C\log_2(D+2)\sqrt D\sum_{i\ge0}2^{-i/2}\le \log_2D+2+3.5\,C\sqrt D\log_2(D+2).
\]
Since $\log_2D+2\le 2\sqrt D\log_2(D+2)$ for $D\ge1$ and $C\ge 1$, we get $u(D)\le 6C\sqrt D\log_2(D+2)$.
\end{proof}

% ====================================================================
\section{Theorem LR: low-rank matrices covered by row and column types}\label{sec:LR}
% ====================================================================

\subsection{Statement}
Let $M\in\R^{R_0\times C_0}$ be a matrix with finite row set $R_0$ and column set $C_0$. A \emph{type} is a function
$\tau$ on $R_0\times C_0$ which is either a \emph{row function}, $\tau(x,w)=a(x)$, or a \emph{column function},
$\tau(x,w)=b(w)$. A constant function is both. A set $\mathcal T$ of types \emph{covers} $M$ on a rectangle
$R\times C$ if for every $(x,w)\in R\times C$ there is $\tau\in\mathcal T$ with $M[x,w]=\tau(x,w)$. For $k,r\ge1$ put
\[
  L_k(r):=2^{14}\,k^3\sqrt{r+2}\,\log_2(r+2).
\]
This is increasing in $k$ and in $r$.

\begin{theorem}[Theorem LR]\label{thm:LR}
Let $M\in\R^{m\times n}$ have rank at most $r$, where $r\ge1$. Let $\mathcal T$ be a set of at most $k$ types
covering $M$ on $[m]\times[n]$. Then there are nonempty $R\subseteq[m]$ and $C\subseteq[n]$ and a single
$\tau^\ast\in\mathcal T$ such that
\[
  |R|\ge 2^{-L_k(r)}m,\qquad |C|\ge 2^{-L_k(r)}n,\qquad M[x,w]=\tau^\ast(x,w)\ \text{ for all }(x,w)\in R\times C .
\]
\end{theorem}

Since $L_k$ is increasing in $k$, the theorem also holds with any $k'\ge k$ in place of $k$. When it is
applied to a family of fewer than $k$ types, the family may be padded by repeating a type.

\subsection{Margin rounding}

\begin{lemma}[margin rounding; quantitative form of Rothvoss~{\cite[Sec.~3]{Rothvoss}}]\label{lem:R}
Let $N\in\R^{m\times n}$ have rank at most $s$, where $s\ge 1$, and suppose $\tfrac12\le|N_{xw}|\le K$ for all entries. Let
$Q^+:=\{(x,w):N_{xw}>0\}$ and $Q^-:=\{(x,w):N_{xw}<0\}$, and suppose $|Q^+|\ge\delta mn$ with
$0<\delta\le\frac12$. Put $T:=\lceil\pi K\sqrt s\,\ln(2/\delta^2)\rceil$. Then there are $X'\subseteq[m]$ and
$Y'\subseteq[n]$ with
\[
|X'|\ge\tfrac\delta2 4^{-T}m,\qquad |Y'|\ge\tfrac\delta24^{-T}n,\qquad |(X'\times Y')\cap Q^-|\le\delta|X'||Y'|.
\]
\end{lemma}
\begin{proof}
Let $s_0:=\rk N$. Then $1\le s_0\le s$, since $N$ has nonzero entries. Write $N_{xw}=\langle u_x,v_w\rangle$ with
$u_x,v_w\in\R^{s_0}$, where the $u_x$ span $\R^{s_0}$. The body $K_0:=\operatorname{conv}\{\pm u_x\}$ is
symmetric with nonempty interior. By Theorem~\ref{thm:John} there is an ellipsoid $E$ with
$E\subseteq K_0\subseteq\sqrt{s_0}E$. Choose a linear map $A$ with $A(E)=s_0^{-1/4}B$, where $B$ is the unit
ball, and replace $u_x$ by $Au_x$ and $v_w$ by $A^{-T}v_w$. This preserves all inner products, and now
$s_0^{-1/4}B\subseteq K_0\subseteq s_0^{1/4}B$. Hence $|u_x|\le s_0^{1/4}$. Moreover $v_w\ne0$, because $N_{xw}\ne0$.
The point $z:=s_0^{-1/4}v_w/|v_w|$ lies in $K_0$, so
\[
s_0^{-1/4}|v_w|=\langle z,v_w\rangle\le\max_x|\langle u_x,v_w\rangle|\le K,
\]
that is, $|v_w|\le Ks_0^{1/4}$. For the unit vectors $\hat u_x,\hat v_w$ we get
$|\langle\hat u_x,\hat v_w\rangle|\ge \frac{1/2}{K\sqrt{s_0}}\ge b:=\frac1{2K\sqrt s}$, with the sign of $N_{xw}$.

Let $g_1,\dots,g_T$ be independent standard Gaussian vectors in $\R^{s_0}$. Put
$X':=\{x:\langle\hat u_x,g_t\rangle\ge0\ \forall t\}$ and $Y':=\{w:\langle\hat v_w,g_t\rangle\ge0\ \forall t\}$, and
$\mathcal R:=X'\times Y'$. By Theorem~\ref{thm:Sheppard} and independence, $\Pr[(x,w)\in\mathcal R]\ge 4^{-T}$ for
$(x,w)\in Q^+$, and $\Pr[(x,w)\in\mathcal R]\le4^{-T}(1-2b/\pi)^T\le 4^{-T}e^{-2bT/\pi}$ for $(x,w)\in Q^-$. Since
$2bT/\pi=T/(\pi K\sqrt s)\ge\ln(2/\delta^2)$, we get
\[
\mathbb E|\mathcal R\cap Q^-|\le\tfrac{\delta^2}{2}mn4^{-T}\le\tfrac\delta2\,\mathbb E|\mathcal R\cap Q^+|,
\qquad \mathbb E|\mathcal R\cap Q^+|\ge\delta mn4^{-T}.
\]
Hence $\mathbb E\bigl[|\mathcal R\cap Q^+|-\delta^{-1}|\mathcal R\cap Q^-|\bigr]\ge\frac12\mathbb E|\mathcal R\cap Q^+|
\ge\frac\delta2mn4^{-T}$. Fix an outcome attaining at least this value. For it,
$|\mathcal R\cap Q^-|\le\delta|\mathcal R\cap Q^+|\le\delta|\mathcal R|$ and
$|\mathcal R|\ge\frac\delta2 mn4^{-T}$. Since $|X'|\le m$ and $|Y'|\le n$, both size bounds follow.
\end{proof}

\subsection{Refinements and the normal form}
A \emph{refinement of cost $(c_1,c_2)$} of a rectangle $R_0\times C_0$ is a rectangle $R\times C\subseteq R_0\times C_0$
with $|R|\ge|R_0|/c_1$ and $|C|\ge|C_0|/c_2$. Every property stated ``everywhere on'' a rectangle is inherited by
its refinements. A \emph{basic quantity} is a positive row function, a positive column function or a positive
constant. A function $\Delta$ on a rectangle is \emph{$c$-comparable} to a basic quantity $Q$ if
$Q\le|\Delta|\le cQ$ everywhere on it.

\begin{lemma}[Lemma A]\label{lem:A}
Let $a$ be a row function and $b$ a column function on $R_0\times C_0$. There is a refinement of cost $(4,2)$ on
which either (i) $a$ is constant, or (ii) for some $\theta$ we have $a(x)>\theta\ge b(w)$ everywhere, or
$a(x)<\theta\le b(w)$ everywhere.
\end{lemma}
\begin{proof}
Let $\theta$ be a median of $a$ on $R_0$, so that at least half of the rows have $a\le\theta$ and at least
half have $a\ge\theta$. If $|\{a=\theta\}|\ge|R_0|/4$, take $\{a=\theta\}\times C_0$; this is case (i).
Otherwise $\{a>\theta\}$ and $\{a<\theta\}$ each contain at least $|R_0|/4$ rows. One of $\{b\le\theta\}$ and
$\{b\ge\theta\}$ contains at least $|C_0|/2$ columns. Take $\{a>\theta\}\times\{b\le\theta\}$ or
$\{a<\theta\}\times\{b\ge\theta\}$ accordingly; this is case (ii).
\end{proof}

\begin{lemma}[Lemma B]\label{lem:B}
Let $\Delta>0$ be a row function and $E\ge0$ a column function on $R_0\times C_0$. There is a refinement of cost
$(4,4)$ on which (a) $\Delta\ge2E$ everywhere, or (b) $E\ge2\Delta$ everywhere, or (c) for some $s>0$ both
$\Delta$ and $E$ take values in $(s/4,4s)$.
\end{lemma}
\begin{proof}
Let $s$ be a median of $E$ on $C_0$. If $s=0$, then $\{E=0\}$ contains at least half of the columns and
$\Delta>0=2E$ there: case (a). Let $s>0$. If $|\{\Delta\ge2s\}|\ge|R_0|/4$, take
$\{\Delta\ge2s\}\times\{E\le s\}$: case (a). Else, if $|\{\Delta\le s/2\}|\ge|R_0|/4$, take
$\{\Delta\le s/2\}\times\{E\ge s\}$: case (b). Otherwise $R_1:=\{s/2<\Delta<2s\}$ contains more than
$|R_0|/2$ rows. If $|\{E\le s/4\}|\ge|C_0|/4$, take $R_1\times\{E\le s/4\}$: case (a). Else, if
$|\{E\ge4s\}|\ge|C_0|/4$, take $R_1\times\{E\ge4s\}$: case (b). Otherwise $\{s/4<E<4s\}$ contains more than
$|C_0|/2$ columns, and $R_1\times\{s/4<E<4s\}$ gives case (c).
\end{proof}

\begin{lemma}[normal form, Lemma NF]\label{lem:NF}
Let a set $\mathcal T$ of $k$ types cover $M$ on $R_0\times C_0$. There is a refinement $R\times C$ of cost at
most $(2^{2k^2+3},2^{2k^2+3})$, and a subset $\mathcal T'\subseteq\mathcal T$ with $k':=|\mathcal T'|\ge1$
that still covers $M$ on $R\times C$, such that on $R\times C$:
\begin{enumerate}[label=(NF\arabic*)]
\item the types of $\mathcal T'$ are strictly ordered, $\tau_{(1)}(x,w)<\dots<\tau_{(k')}(x,w)$, in the same order at
  every $(x,w)$;
\item if $k'\ge2$, every gap $g_j:=\tau_{(j+1)}-\tau_{(j)}$ is $16$-comparable to a basic quantity $Q_j$;
\item if $k'\ge2$, there is $j^\ast$ with $Q_{j^\ast}\ge Q_j/16$ for all $j$.
\end{enumerate}
\end{lemma}
\begin{proof}
Process the $\binom k2$ pairs $\{\tau,\tau'\}$ one after the other, refining each time. If one of the two types
has already been deleted, skip the pair. A constant is treated as a row function when it is paired with a row
function, and as a column function when it is paired with a column function (editorial item~(i) of the referee report on Theorem~LR).
\begin{itemize}[leftmargin=1.5em]
\item \emph{Two row functions} (or a row function and a constant, or two constants). Split the rows into three
  parts according to the sign of $\tau-\tau'$, and keep a largest part: cost $(3,1)$. If the part with
  $\tau=\tau'$ is kept, delete $\tau'$; coverage is preserved. Otherwise $|\tau-\tau'|$ is a positive row
  function, hence $1$-comparable to a basic quantity, and has constant sign. Two column functions are treated
  symmetrically, at cost $(1,3)$.
\item \emph{A row function $a$ and a column function $b$} (neither constant). Apply Lemma~\ref{lem:A}, at cost
  $(4,2)$. In case~(i) $a$ is now constant, and from now on it is treated as a column function. Split the
  columns by the sign of $a-b$ (cost $(1,3)$), deleting $b$ if the part with $a=b$ is kept. In case~(ii), say
  $a>\theta\ge b$ (the other case is symmetric), we have $a-b=\Delta+E$ with $\Delta:=a-\theta>0$ a row
  function and $E:=\theta-b\ge0$ a column function. Lemma~\ref{lem:B}, at cost $(4,4)$, gives one of three
  cases: (a) $\Delta\le a-b\le\frac32\Delta$; (b) $E\le a-b\le\frac32E$, where $E\ge2\Delta>0$; or
  (c) $s/2<a-b<8s=16\cdot(s/2)$. In each case $a-b>0$ is $16$-comparable to a basic quantity.
\end{itemize}
Each pair costs at most $(16,8)$, so all pairs together cost at most $16^{\binom k2}$ per side. On the
resulting rectangle, every difference of two surviving types has a constant strict sign and is $16$-comparable to
a basic quantity. Since the values at a fixed $(x,w)$ are real numbers, these signs define one strict total
order, and every entry has exactly one type in $\mathcal T'$. This gives (NF1) and~(NF2).

For (NF3), let $J_{\mathrm{row}}$, $J_{\mathrm{col}}$, $J_{\mathrm{const}}$ be the sets of $j$ for which $Q_j$ is a row function,
a column function or a constant. Split the rows according to which $j\in J_{\mathrm{row}}$ maximises $Q_j(x)$, keeping
a largest part (cost at most $k$). On the result some $j_r$ satisfies $Q_{j_r}=R^\ast:=\max_{J_{\mathrm{row}}}Q_j$.
Do the same for the columns, which gives $j_c$ with $Q_{j_c}=C^\ast:=\max_{J_{\mathrm{col}}}Q_j$. Let $j_s$ attain
$s^\ast:=\max_{J_{\mathrm{const}}}Q_j$. Apply Lemma~\ref{lem:B} to $(R^\ast,C^\ast)$.
In case (a) ($R^\ast\ge2C^\ast$), split the rows into $\{R^\ast\ge s^\ast\}$ and $\{R^\ast<s^\ast\}$ (cost $2$).
On the first part $j^\ast=j_r$ works, and on the second $j^\ast=j_s$ works. Case (b) is symmetric, splitting the
columns. In case (c) ($R^\ast,C^\ast\in(s/4,4s)$): if $s^\ast\ge s$, then $j^\ast=j_s$ works with factor~$4$;
if $s^\ast<s$, then $R^\ast>s/4>s^\ast/4$ and $C^\ast<4s<16R^\ast$, so $j^\ast=j_r$ works with factor~$16$.
If some $J_\bullet$ is empty, the corresponding comparison is simply omitted. The cost of this step is at most
$(8k,8k)$. The total cost is at most $16^{\binom k2}\cdot 8k\le2^{2k^2+3}$ per side.
\end{proof}

\subsection{Proof of Theorem LR}
We use induction on $k$. For $k=1$ take $R=[m]$ and $C=[n]$. Let $k\ge2$, and apply Lemma~\ref{lem:NF} to
$[m]\times[n]$, obtaining $R_0\times C_0$ and $\mathcal T'$. If $k'=1$ we are done. Otherwise, by (NF2) and
(NF3), $g_{j^\ast}\ge Q_{j^\ast}$ and $g_j\le16Q_j\le256Q_{j^\ast}$ for every $j$. So the spread
$\mathcal S:=\tau_{(k')}-\tau_{(1)}=\sum_jg_j$ satisfies $\mathcal S\le256(k'-1)Q_{j^\ast}$. Put
\[
z:=\tfrac12\bigl(\tau_{(j^\ast)}+\tau_{(j^\ast+1)}\bigr),\qquad \sigma:=Q_{j^\ast},\qquad N:=(M-z)/\sigma
\quad\text{(entrywise, on $R_0\times C_0$)}.
\]
Each $\tau$ is a row or a column function, so $z=\varphi(x)+\vartheta(w)$ for functions $\varphi,\vartheta$.
Since $\sigma$ is a basic quantity, $1/\sigma=\alpha(x)\beta(w)$ with $\alpha,\beta>0$ (one of them $\equiv1$).
Hence $N=\mathrm{diag}(\alpha)\,(M-\varphi\mathbf 1^T-\mathbf 1\vartheta^T)\,\mathrm{diag}(\beta)$ has rank at
most $r+2$. By (NF1) each entry has a unique type $\tau_{(i)}$. If $i\le j^\ast$, then
$-256(k'-1)\le-\mathcal S/\sigma\le N\le-g_{j^\ast}/(2\sigma)\le-\frac12$. If $i>j^\ast$, then
$\frac12\le N\le256(k'-1)$. So $\frac12\le|N|\le K:=256(k-1)$, and the sign of $N$ records on which side of the
gap $j^\ast$ the type of the entry lies.

Replacing $N$ by $-N$ if necessary, we may assume that $Q^+=\{N>0\}$ contains at least half of the entries of
$R_0\times C_0$. Let $\mathcal T_{\mathrm{maj}}$ be the types on the positive side and $\mathcal T_{\mathrm{min}}$ those on the
negative side. Both are nonempty, because $1\le j^\ast<k'$, so $|\mathcal T_{\mathrm{maj}}|\le k'-1\le k-1$. If
$Q^-=\emptyset$, every entry of $R_0\times C_0$ has a type in $\mathcal T_{\mathrm{maj}}$, and we apply the induction
hypothesis to $\mathcal T_{\mathrm{maj}}$ on $R_0\times C_0$ directly. Otherwise put $\delta:=1/(16(r+1)^{k-1})$ and
apply Lemma~\ref{lem:R} (with $s=r+2$). This gives $X'\times Y'\subseteq R_0\times C_0$ with
$|X'|\ge\frac\delta24^{-T}|R_0|$ and $|Y'|\ge\frac\delta24^{-T}|C_0|$, containing at most $\delta|X'||Y'|$ entries of
$Q^-$, where $T=\lceil\pi K\sqrt{r+2}\ln(2/\delta^2)\rceil$.

The entries of $Q^-$ are \emph{exactly} the support of
\[
  H:=\mathop{\bigcirc}_{\tau\in\mathcal T_{\mathrm{maj}}}(M-\tau).
\]
Indeed, an entry whose type lies in $\mathcal T_{\mathrm{maj}}$ makes one factor vanish. By (NF1), an entry whose
type lies in $\mathcal T_{\mathrm{min}}$ differs from every majority-type value at that position. Each matrix
$M-\tau$ has the form $M-a\mathbf 1^T$ or $M-\mathbf 1b^T$, of rank at most $r+1$. By Lemma~\ref{lem:HP},
$\rk H\le(r+1)^{k-1}$. So $H$ restricted to $X'\times Y'$ has at most $|X'||Y'|/(16(r+1)^{k-1})$ nonzero entries,
and Lemma~\ref{lem:GL} gives $X''\times Y''$ with $|X''|\ge|X'|/2$, $|Y''|\ge|Y'|/2$ and $H=0$ there. Every entry
of $X''\times Y''$ therefore has its type in $\mathcal T_{\mathrm{maj}}$. The submatrix has rank at most $r$, and we
apply the induction hypothesis to it with the at most $k-1$ types $\mathcal T_{\mathrm{maj}}$.

\emph{Loss.} We bound the loss of one round, measured in bits per side, i.e.\ $\log_2$ of the ratio of side
lengths. Lemma~\ref{lem:NF} costs at most $2k^2+3$, Lemma~\ref{lem:GL} costs $1$, and Lemma~\ref{lem:R} costs
$\log_2(2/\delta)+2T$. Here $\log_2(2/\delta)=5+(k-1)\log_2(r+1)$. Since $r\ge1$ we have $\ln(r+2)\ge1$, so
$\ln(2/\delta^2)=\ln 512+2(k-1)\ln(r+1)\le 6k\ln(r+2)$. Hence
$T\le1+256\pi(k-1)\sqrt{r+2}\cdot 6k\ln(r+2)\le 1+3345\,k^2\sqrt{r+2}\log_2(r+2)$. Using
$\sqrt{r+2}\log_2(r+2)\ge1$, one round costs at most
\[
2k^2+11+k\log_2(r+2)+6690\,k^2\sqrt{r+2}\log_2(r+2)\le 2^{13}k^2\sqrt{r+2}\log_2(r+2)
\]
bits per side. The number of types drops by at least one per round, so the total loss is at most
$\sum_{j=2}^k2^{13}j^2\sqrt{r+2}\log_2(r+2)\le L_k(r)$. \qed

\subsection{The column-type special case}

\begin{corollary}[RQ-list]\label{cor:RQlist}
Let $M\in\R^{m\times n}$ have rank at most $r$, where $r\ge 1$, and suppose every column of $M$ takes at most
$k$ distinct values. Then there are $R\subseteq[m]$ and $C\subseteq[n]$ with $|R|\ge2^{-L_k(r)}m$ and
$|C|\ge2^{-L_k(r)}n$ such that every column of $M$ is constant on $R\times C$.
\end{corollary}
\begin{proof}
List the values of column $w$ as $v_1(w)<\dots<v_{k_w}(w)$, and pad the list to length $k$ by repeating
$v_{k_w}(w)$. The $k$ column functions $v_i$ cover $M$. Theorem~\ref{thm:LR} gives a rectangle on which
$M[x,w]=v_i(w)$ for a single $i$.
\end{proof}
See Remark~\ref{rem:SS} for the relation of Corollary~\ref{cor:RQlist} to~\cite{SingerSudan}. In the geometric
applications, Theorem~\ref{thm:LR} is always used together with the following observation.

\begin{lemma}[types give orthogonal splits]\label{lem:typeOR}
Let $X\subset\R^D$ be rainbow-$n$-free, and let $R,C\subseteq X$ be disjoint and nonempty. Suppose that on
$R\times C$ the matrix $[\dd xa^2]_{x\in R,a\in C}$ equals a single type $\tau$ of one of the following forms:
a constant $v^2$; a column function $v(a)^2$; or a row function $\dd xs^2$ with $s\in X\setminus C$ fixed. Then
$\min(|R|,|C|)\le h$.
\end{lemma}
\begin{proof}
In the first two cases every $a\in C$ is equidistant from $R$, so Lemma~\ref{lem:OR}(a) applies. In the third,
$\dd xa=\dd xs$ means $x\in\Bis(s,a)$ for all $x\in R$ and $a\in C$, so Lemma~\ref{lem:OR}(b) applies with
$p=s$.
\end{proof}

% ====================================================================
\section{The label lemma and pairs with disjoint labels (G1)}\label{sec:G1}
% ====================================================================

Throughout this section $n\ge4$, $X\subset\R^D$ is finite and rainbow-$n$-free, and $h=s_n(\lfloor D/2\rfloor)$.
Let $T\subseteq X$ be a rainbow set with $|T|=n-2$. For $x\notin T$ the \emph{label set} of $x$ is
$L_x:=\{\dd xt:t\in T\}$. A point $x$ is \emph{$T$-generic} if $x\notin T$ and $T\cup\{x\}$ is rainbow; then
$|L_x|=n-2$. Let $Y_T$ be the set of $T$-generic points of $X$. Two points $x,w$ are \emph{label-disjoint} if
$L_x\cap L_w=\emptyset$. Put
\[
k_G:=\binom{n-2}{2}+2(n-2),\qquad \mu_G:=2^{L_{k_G}(D+2)}(h+1).
\]

\begin{lemma}[label lemma]\label{lem:label}
If $x,w\in Y_T$ are label-disjoint, then $x\ne w$ and $\dd xw\in\Dist(T)\cup L_x\cup L_w$.
\end{lemma}
\begin{proof}
$L_x\ne\emptyset$, so $x\ne w$. The $n$-set $T\cup\{x,w\}$ is not rainbow. Its distances other than $\dd xw$
are $\Dist(T)\cup L_x\cup L_w$, and these $\binom n2-1$ numbers are pairwise distinct: $\Dist(T)\cup L_x$ and
$\Dist(T)\cup L_w$ are rainbow, and $L_x\cap L_w=\emptyset$. Hence $\dd xw$ equals one of them.
\end{proof}

\begin{theorem}[G1, bipartite form]\label{thm:G1bip}
Let $P,Q\subseteq Y_T$ be such that every $x\in P$ and $w\in Q$ are label-disjoint. Then
$\min(|P|,|Q|)\le 2^{L_{k_G}(D+2)}h\le\mu_G$.
\end{theorem}
\begin{proof}
We may assume $P,Q\neq\emptyset$. Then $P\cap Q=\emptyset$, by Lemma~\ref{lem:label}. By Lemma~\ref{lem:label},
for $x\in P$ and $w\in Q$ the entry $\dd xw^2$ of the matrix $[\dd xw^2]_{x\in P,w\in Q}$ equals one of the following
$k_G$ types: the constants $d^2$ with $d\in\Dist(T)$; the column functions $w\mapsto\dd wt^2$ with $t\in T$; and the
row functions $x\mapsto\dd xt^2$ with $t\in T$. The matrix has rank at most $D+2$ (Lemma~\ref{lem:distrank}).
Theorem~\ref{thm:LR} gives $R\subseteq P$ and $C\subseteq Q$ with $|R|\ge2^{-L}|P|$ and $|C|\ge2^{-L}|Q|$, where
$L=L_{k_G}(D+2)$, and a single type on $R\times C$. The points of $T$ are not in $Q\subseteq Y_T$, so
Lemma~\ref{lem:typeOR} applies and gives $\min(|R|,|C|)\le h$. Hence $\min(|P|,|Q|)\le2^Lh$.
\end{proof}

\begin{corollary}[G1, counting forms]\label{cor:G1}
Let $\mu\ge\mu_G$.
\begin{enumerate}[label=(\alph*)]
\item For every $Y\subseteq Y_T$, the number of ordered pairs $(x,w)\in Y^2$ that are label-disjoint is at most
  $4^{n-2}\mu|Y|$.
\item For disjoint $P,Q\subseteq Y_T$, the number of label-disjoint pairs in $P\times Q$ is at most
  $4^{n-2}\mu(|P|+|Q|)$.
\item (degree form) For $Y\subseteq Y_T$ and a real $v$ put $\deg_Y(v):=|\{x\in Y:v\in L_x\}|$. If
  $\deg_Y(v)\le|Y|/(4(n-2))$ for every $v$, then $|Y|\le\frac43(4^{n-2}\mu+1)\le2\cdot4^{n-2}\mu+4$.
\end{enumerate}
\end{corollary}
\begin{proof}
(a) Let $N_{\mathrm{ld}}$ be the number of label-disjoint ordered pairs. Colour every element of
$\bigcup_{x\in Y}L_x$ red or blue, independently and uniformly. Let $P$ be the set of $x\in Y$ all of whose labels
are red, and $Q$ the set of $w\in Y$ all of whose labels are blue. Every pair in $P\times Q$ is label-disjoint, so
by Theorem~\ref{thm:G1bip}, $|P||Q|\le\mu\max(|P|,|Q|)\le\mu|Y|$. A label-disjoint pair $(x,w)$ involves
$2(n-2)$ distinct labels, so it lies in $P\times Q$ with probability $4^{-(n-2)}$. Taking expectations,
$4^{-(n-2)}N_{\mathrm{ld}}\le\mathbb E|P||Q|\le\mu|Y|$.
(b) Apply (a) to $Y=P\cup Q$.
(c) The ordered pairs $(x,w)$ with $x\ne w$ that share a label number at most
$\sum_v\deg_Y(v)^2\le\max_v\deg_Y(v)\cdot\sum_v\deg_Y(v)\le\frac{|Y|}{4(n-2)}(n-2)|Y|=|Y|^2/4$. So
$N_{\mathrm{ld}}\ge|Y|(|Y|-1)-|Y|^2/4$, and (a) gives $\frac34|Y|^2-|Y|\le4^{n-2}\mu|Y|$.
\end{proof}

In the sequel G1 is applied with $\mu:=2^{\Lst}(h+1)$, where $\Lst=L_{\kst}(D+2)$ and $\kst\ge k_G$
(Section~\ref{sec:main}). Then $\mu\ge\mu_G$, since $L_k$ is increasing in $k$.

% ====================================================================
\section{Configurations, the relative bound HR and the kill step}\label{sec:framework}
% ====================================================================

\subsection{Standing assumptions and absolute parameters}
From now on $n\ge4$ is fixed, $D\ge1$, $X\subset\R^D$ is a finite rainbow-$n$-free set, and
$h:=s_n(\lfloor D/2\rfloor)$. All parameters below are explicit. Their roles are explained where they are used,
but they are defined here, so that every later definition refers only to earlier ones.

\begin{definition}[type counts]\label{def:typecounts}
Put $t:=\binom{n-2}2$ and $k':=t+2n$. For $0\le r\le n-5$ put $p_r:=n-4-r$ and
\begin{gather*}
w_1(r):=(n-2)+t+p_r,\quad w_2(r):=t+p_r,\quad k_3(w):=w+1+r,\quad k_4(w):=w+2r+\tbinom r2,\\
k_5:=\tbinom{n-4}2+2(n-4),\qquad k_c(r):=k_3(w_1(r))+k_4(w_1(r)).
\end{gather*}
For $0\le r\le n-5$ and $1\le p\le n-4-r$ put $K^{(r,p)}:=12\cdot6^{\,n-4-r-p}$, and let
$K_r:=K^{(r,1)}=12\cdot6^{\,n-5-r}$. For $0\le r\le n-6$ and $1\le p\le n-5-r$ put
\[
k_{\mathrm{sl}}(r,p):=3K^{(r,p+1)}-1+\tbinom{r+p}2+2p+2r .
\]
Let $\kst$ be the maximum of $k'$, of all $k_c(r)$ and of all $k_{\mathrm{sl}}(r,p)$. For $n=4$ only $k'$ is
defined, and $\kst:=k'$.
\end{definition}

One checks directly that $k_G\le k'\le\kst$, and that $k_5$, $k_3(w_2)$, $k_4(w_2)$, $k_3(w_1)$, $k_4(w_1)$ are all at
most $k_c(r)\le\kst$. For instance $k_3(w_1)=\binom{n-4}2+4n-12\ge k_5$, using $r+p_r=n-4$. Also
$\kst\le 6^{\,n-4}+2n^2$.

\begin{definition}[absolute parameters]\label{def:absolute}
\[
\Lst:=L_{\kst}(D+2),\quad \mu:=2^{\Lst}(h+1),\quad e:=2^{\Lst+1}(h+1),\quad e_A:=2^{\Lst+2}(h+1),\quad
\smin:=2^{\Lst+2+2n}(h+1).
\]
\end{definition}
These depend only on $n$, $D$ and $h$. We shall use the following inequalities, all immediate from the
definitions (recall $h\ge1$, so $\mu\ge2$):
\begin{equation}\label{eq:absorb}
1\le\mu\le e\le e_A\le\smin,\qquad 4^n\mu\le\smin,\qquad 2\cdot4^{n-2}\mu+4\le\smin,\qquad
e>2^{\Lst+1}h .
\end{equation}
Since $\Lst\ge L_{k}(D+2)$ for every $k\le\kst$, any rectangle given by Theorem~\ref{thm:LR} with at most $\kst$ types
for a matrix of rank at most $D+2$ has both sides of relative size at least $2^{-\Lst}$.

\subsection{Configurations}

\begin{definition}[level-$j$ configuration]\label{def:config}
Let $0\le j\le n-2$. A \emph{level-$j$ configuration} is a triple $(S,W,\sigma)$ consisting of a rainbow set
$S\subseteq X$ with $|S|=j$, a set $W\subseteq X\setminus S$ such that $S\cup\{x\}$ is rainbow for every $x\in W$, and a
real $\sigma>0$ (the \emph{cap}) such that
\[
|W\cap\Sph(s,t)|<\sigma\qquad\text{for all }s\in S,\ t>0 .
\]
For $j=0$ the cap condition is void, and $(\emptyset,W,\sigma)$ is a configuration for every $W\subseteq X$
and every $\sigma>0$. For $b\in W$ put
\[
R_b:=\Dist(S\cup\{b\})=\Dist(S)\cup\{\dd bs:s\in S\},\qquad
\Bg_b(\tau):=\{\lambda>0:\lambda\notin R_b,\ |W\cap\Sph(b,\lambda)|\ge\tau\}\quad(\tau>0).
\]
\end{definition}

If $(S,W,\sigma)$ is a level-$j$ configuration, $W'\subseteq W$ and $\sigma'\ge\sigma$, then $(S,W',\sigma')$ is
again a level-$j$ configuration.

\begin{definition}[Claim$(j)$]\label{def:claim}
For $0\le j\le n-2$ and a real $\Gamma_j\ge1$, \textbf{Claim$(j)$} (with constant $\Gamma_j$) is the statement:
\emph{every level-$j$ configuration $(S,W,\sigma)$ with $\sigma\ge\smin$ satisfies $|W|\le\Gamma_j\sigma$.}
\end{definition}

Claim$(j)$ quantifies over \emph{all} configurations and \emph{all} caps $\sigma\ge\smin$, and $\Gamma_j$ does
not depend on $\sigma$. This uniformity is what allows the lemmas below to apply Claim$(j+1)$ with caps
other than $\sigma$.

\begin{lemma}[extension]\label{lem:ext}
Let $(S,W,\sigma)$ be a level-$j$ configuration, and let $b\in W$ and $x\in W\setminus\{b\}$. Suppose that
$\dd xb\notin R_b$, that $x\notin\bigcup_{s\in S}\Bis(b,s)$, and that $\dd xs\ne\dd b{s'}$ for all $s,s'\in S$.
Then $S\cup\{b,x\}$ is rainbow. Moreover, at most $j^2\sigma$ points $x\in W$ satisfy $\dd xs=\dd b{s'}$ for some
$s,s'\in S$; there are no such points if $j=0$.
\end{lemma}
\begin{proof}
The distances of $S\cup\{b,x\}$ are those of $S$, the numbers $\dd bs$ and $\dd xs$ ($s\in S$), and $\dd bx$. The
sets $\Dist(S)\cup\{\dd bs\}$ and $\Dist(S)\cup\{\dd xs\}$ are rainbow because $b,x\in W$. By hypothesis
$\dd bs\ne\dd x{s'}$, $\dd bx\notin\Dist(S)\cup\{\dd bs\}$ and $\dd bx\ne\dd xs$, and $b,x\notin S$ are distinct.
For the second statement: such points lie on one of the $j^2$ spheres $\Sph(s,\dd b{s'})$, each of which meets
$W$ in fewer than $\sigma$ points.
\end{proof}

\subsection{The hereditary relative bound}

\begin{lemma}[HR$_r$]\label{lem:HR}
Let $0\le r\le n-3$ and assume Claim$(r+1)$ with constant $\Gamma_{r+1}$. Put $\Gamma'_r:=\Gamma_{r+1}+r^2$.
Let $(S,W,\sigma)$ be a level-$r$ configuration with $\sigma\ge\smin$, let $c\in W$, $V\subseteq W\setminus\{c\}$ and
$\theta\ge\sigma$. Then
\[
\Bigl|\Bigl\{x\in V:\ \dd xc\notin R_c,\ x\notin\textstyle\bigcup_{s\in S}\Bis(c,s),\ |V\cap\Sph(c,\dd xc)|<\theta\Bigr\}\Bigr|
\;\le\;\Gamma'_r\,\theta .
\]
\end{lemma}
\begin{proof}
Let $V_0$ be the set on the left. Let $V_2$ be the set of $x\in V_0$ with $\dd xs=\dd c{s'}$ for some
$s,s'\in S$. By Lemma~\ref{lem:ext}, $|V_2|\le r^2\sigma\le r^2\theta$. Put $V_1:=V_0\setminus V_2$. For
$x\in V_1$ the set $S\cup\{c,x\}$ is rainbow, by Lemma~\ref{lem:ext}. We claim that $(S\cup\{c\},V_1,\theta)$ is a
level-$(r+1)$ configuration. Indeed, $S\cup\{c\}$ is rainbow because $c\in W$, and $V_1\subseteq X\setminus(S\cup\{c\})$.
A sphere about $s\in S$ meets $V_1\subseteq W$ in fewer than $\sigma\le\theta$ points. If a sphere $\Sph(c,t)$
meets $V_1$ in a point $x$, then $t=\dd xc$ and $|V_1\cap\Sph(c,t)|\le|V\cap\Sph(c,\dd xc)|<\theta$. Since
$\theta\ge\smin$, Claim$(r+1)$ gives $|V_1|\le\Gamma_{r+1}\theta$.
\end{proof}

\begin{remark}[the bisector exception]\label{rem:bisector}
Points of $V$ on the bisectors $\Bis(c,s)$ are not controlled by Lemma~\ref{lem:HR}, and neither are points on
spheres about $c$ with radius in $R_c$. In every application such points are handled by Theorem~\ref{thm:LR}. A
point $x\in\Bis(c,s)$ satisfies $\dd xc^2=\dd xs^2$, which is a \emph{row type} of the matrix
$[\dd xc^2]$ (rows $x$, columns $c$). A single-type rectangle of this type has $R\subseteq\Bis(s,c)$ for all its
columns $c$, so Lemma~\ref{lem:typeOR} applies. The radii in $R_c$ are column functions ($\dd cs$) or constants
($d\in\Dist(S)$).
\end{remark}

\subsection{The kill step}

\begin{lemma}[kill step]\label{lem:kill}
Let $0\le r\le n-3$ and assume Claim$(r+1)$ with constant $\Gamma_{r+1}$. Let $(S,W,\sigma)$ be a level-$r$
configuration with $\sigma\ge\smin$, let $\tau\ge\sigma$, and let $K\ge0$ be an integer such that
$|\Bg_b(\tau)|\le K$ for every $b\in W$. Then
\[
|W|\le P(K,r)\,\bigl(2\Gamma_{r+1}\tau+2r^2\sigma+1\bigr),\qquad P(K,r):=\binom{D+1+e_{K,r}}{e_{K,r}}^{2},\quad
e_{K,r}:=K+2r+\tbinom r2 .
\]
\end{lemma}
\begin{proof}
For $b\in W$ define
\[
\Delta_b(y):=\prod_{s\in S}\bigl(\dd yb^2-\dd ys^2\bigr)\prod_{s\in S}\bigl(\dd yb^2-\dd bs^2\bigr)
\prod_{d\in\Dist(S)}\bigl(\dd yb^2-d^2\bigr),\qquad
F_b(y):=\Delta_b(y)\prod_{\lambda\in\Bg_b(\tau)}\bigl(\dd yb^2-\lambda^2\bigr).
\]
$F_b$ is a product of at most $e_{K,r}$ affine functions of the lifted coordinates, so by Lemma~\ref{lem:DIM}
all $F_b$ lie in one space of dimension $\binom{D+1+e_{K,r}}{e_{K,r}}$. At $y=b$ the factors equal $-\dd bs^2$,
$-\dd bs^2$, $-d^2$ and $-\lambda^2$, all nonzero (recall $b\notin S$), so $F_b(b)\ne0$.

\emph{Claim: for every $b\in W$, fewer than $\Gamma_{r+1}\tau+r^2\sigma$ points $x\in W\setminus\{b\}$ have
$F_b(x)\ne0$.} Let $F_b(x)\ne0$. Then $\dd xb\notin\Bg_b(\tau)\cup R_b$ and $x\notin\bigcup_s\Bis(b,s)$. By
Lemma~\ref{lem:ext}, either $x$ is one of at most $r^2\sigma$ points with $\dd xs=\dd b{s'}$, or $S\cup\{b,x\}$ is
rainbow and $x\in W':=\{x\in W\setminus\{b\}:S\cup\{b,x\}\text{ rainbow},\ \dd xb\notin\Bg_b(\tau)\}$. The triple
$(S\cup\{b\},W',\tau)$ is a level-$(r+1)$ configuration. Indeed, $S\cup\{b\}$ is rainbow. Spheres about $s\in S$
meet $W'$ in fewer than $\sigma\le\tau$ points. If $x\in W'\cap\Sph(b,t)$, then $t=\dd xb\notin R_b$, because
$S\cup\{b,x\}$ is rainbow, and $t\notin\Bg_b(\tau)$; hence $|W'\cap\Sph(b,t)|\le|W\cap\Sph(b,t)|<\tau$. Since
$\tau\ge\smin$, Claim$(r+1)$ gives $|W'|\le\Gamma_{r+1}\tau$.

Let $G$ be the graph on $W$ in which $\{x,b\}$ is an edge if $F_b(x)\ne0$ and $F_x(b)\ne0$. By the claim,
$G$ has at most $|W|(\Gamma_{r+1}\tau+r^2\sigma)$ edges, so its average degree is at most
$2(\Gamma_{r+1}\tau+r^2\sigma)$. By Theorem~\ref{thm:CW}, $G$ has an independent set $I$ with
$|I|\ge|W|/(2\Gamma_{r+1}\tau+2r^2\sigma+1)$. For distinct $x,y\in I$ we have $F_x(y)=0$ or $F_y(x)=0$, so
Lemma~\ref{lem:H} gives $|I|\le\binom{D+1+e_{K,r}}{e_{K,r}}^2=P(K,r)$.
\end{proof}

\begin{remark}[on the constants of Lemma~\ref{lem:kill}]\label{rem:killconst}
The square in $P(K,r)$ comes from Lemma~H, where a pair is killed from one side only. The term $+1$ comes from
Caro--Wei. When the lemma is applied with $\tau=\rho\sigma$ for a number $\rho\ge1$, we use $1\le\sigma$
(by~\eqref{eq:absorb}) to write the bound as $|W|\le P(K,r)(2\rho\Gamma_{r+1}+2r^2+1)\,\sigma$. This is a pure
multiple of the cap.
\end{remark}

% ====================================================================
\section{The top levels: Theorem R$^+$ with a single cap}\label{sec:top}
% ====================================================================

In this section we prove Claim$(j)$ for $j=n-2,n-3,n-4$ (for $n=4$: $j=2,1,0$). At these levels a point
with two big classes is handled directly, using $n-4$ points of $S$ as ``spread'' points; no class-count
statement is needed.

\begin{theorem}[Theorem R$^+$]\label{thm:Rplus}
Put $Q:=2^{44}n^4(D+2)^{8k'}$, $\Gamma_{n-2}:=4(n-2)^2$, and for $j=n-3,n-4$
\begin{equation}\label{eq:Gamma-top}
\Gamma_j:=P^+_j\,\bigl(2Q\Gamma_{j+1}+2(j^2+2n)+1\bigr),\qquad P^+_j:=\binom{D+1+e_j}{e_j}^{2},\quad
e_j:=2+2j+\tbinom j2 .
\end{equation}
Then Claim$(j)$ holds with constant $\Gamma_j$ for $j=n-2,n-3,n-4$.
\end{theorem}

\begin{remark}[provenance and referee fixes]\label{rem:Rplus-prov}
This is Theorem~R$^+$ of the project notes. There it was formulated with a sequence of thresholds
$\tau_0<\tau_1<\cdots$, one for each point of $S$. An AI referee pass (report ``REFEREE\_R'') found it correct
subject to three fixes: \textbf{F1}, a genuine threshold gap in Step~A, repaired by the branch at
$\theta=Q^{1/2}\sigma$ in Step~4 below; \textbf{F2}, the auxiliary points $a$ must be taken in $W\setminus\{b\}$;
\textbf{F3}, an undefined threshold $\tau_{-1}$, which is moot here because only the points of $S'$ are used and
$S'=\emptyset$ is allowed. A later referee pass (``referee15'', item~X1) asked that the single-cap version
be written out in full rather than obtained by a translation remark. This section does so. In particular the
two-class count (``Lemma~S2'') and the ``four-case analysis'' of the notes appear below as Step~3 and
Lemma~\ref{lem:twodist}. The ``level switch'' for $n\ge6$ that was sketched in the notes after that referee pass is not
needed: the configuration formalism of Section~\ref{sec:framework} replaces it. For $n=4$ the Main Theorem follows from
Theorem~\ref{thm:Rplus} alone. An earlier, longer proof of the case $n=4$ in the project notes, together with the editorial
items of its referee report that concern only that proof, is superseded and not used here.
\end{remark}

\subsection{The base level}

\begin{proof}[Proof of Claim$(n-2)$]
Let $(S,W,\sigma)$ be a level-$(n-2)$ configuration with $\sigma\ge\smin$. Then $S$ is a rainbow $(n-2)$-set and every
$x\in W$ is $S$-generic, i.e.\ $W\subseteq Y_S$ in the notation of Section~\ref{sec:G1}. For every real $v$,
$\deg_W(v)\le\sum_{s\in S}|W\cap\Sph(s,v)|<(n-2)\sigma$. Suppose $|W|\ge4(n-2)^2\sigma$. Then
$\deg_W(v)<|W|/(4(n-2))$ for all $v$, and Corollary~\ref{cor:G1}(c) gives
$|W|\le2\cdot4^{n-2}\mu+4\le\smin\le\sigma<4(n-2)^2\sigma$, a contradiction. Hence
$|W|<4(n-2)^2\sigma=\Gamma_{n-2}\sigma$.
\end{proof}

\subsection{Two distributions}

\begin{lemma}[two-distribution lemma]\label{lem:twodist}
Let $p=(p_t)$ and $q=(q_t)$ be finitely supported probability distributions on $\R$, let $\rho\ne\rho'$ be reals and
$0\le\eta\le1/16$. Suppose
\[
(1-p_{\rho'})(1-q_\rho)-\sum_{t\notin\{\rho,\rho'\}}p_tq_t\;\le\;\eta .
\]
Then there is a real $t^\ast$ with $p_{\rho'}+p_{t^\ast}\ge1-\eta^{1/4}$ or $q_\rho+q_{t^\ast}\ge1-\eta^{1/4}$.
\end{lemma}
\begin{proof}
Let $\alpha:=1-p_{\rho'}$ and $\beta:=1-q_\rho$. If $\alpha\le\eta^{1/4}$ or $\beta\le\eta^{1/4}$, the claim
holds with any $t^\ast$. Otherwise $\alpha\beta>\sqrt\eta$. For $t\notin\{\rho,\rho'\}$ put $p'_t:=p_t/\alpha$ and
$q'_t:=q_t/\beta$. Then $p'_t,q'_t\le1$, $\sum_tp'_t\le1$ and $\sum_tq'_t\le1$ (sums over $t\notin\{\rho,\rho'\}$),
and the hypothesis gives $\sum_tp'_tq'_t\ge1-\eta/(\alpha\beta)\ge1-\epsilon$ with $\epsilon:=\sqrt\eta<1$. Let
$t^\ast\notin\{\rho,\rho'\}$ maximise $q'_t$. Then $1-\epsilon\le\sum_tp'_tq'_t\le q'_{t^\ast}$, and
$1-\epsilon\le p'_{t^\ast}q'_{t^\ast}+\sum_{t\ne t^\ast}q'_t\le p'_{t^\ast}+(1-q'_{t^\ast})\le p'_{t^\ast}+\epsilon$.
So $p'_{t^\ast}\ge1-2\epsilon$, and $p_{\rho'}+p_{t^\ast}=(1-\alpha)+\alpha p'_{t^\ast}\ge1-2\sqrt\eta\ge1-\eta^{1/4}$,
where the last step uses $\eta\le1/16$.
\end{proof}

\subsection{A column with two big classes (Step A)}

For a level-$j$ configuration $(S,W,\sigma)$ and $b\in W$ let $\Delta_b$ be the polynomial of
Lemma~\ref{lem:kill}:
$\Delta_b(y)=\prod_{s}(\dd yb^2-\dd ys^2)\prod_s(\dd yb^2-\dd bs^2)\prod_{d\in\Dist(S)}(\dd yb^2-d^2)$. If
$\Delta_b(a)\ne0$ then $\dd ab\notin R_b$ and $a\notin\bigcup_s\Bis(b,s)$.

\begin{lemma}[Step A]\label{lem:stepA}
Let $\max(n-4,0)\le j\le n-3$, and let $(S,W,\sigma)$ be a level-$j$ configuration with $\sigma\ge\smin$. Let $b\in W$
and let $\rho_1\ne\rho_2$ be elements of $\Bg_b(Q\sigma)$. Put $C_\nu:=W\cap\Sph(b,\rho_\nu)$ for $\nu=1,2$, so that
$|C_\nu|\ge Q\sigma$. Then the set
\[
A:=\{a\in W\setminus\{b\}:\ \dd ab\notin\{\rho_1,\rho_2\},\ \Delta_b(a)\ne0\}
\]
has $|A|\le e_A+(j^2+2(n-4))\sigma$.
\end{lemma}

\begin{proof}
Fix $S'\subseteq S$ with $|S'|=n-4$; this is possible since $j\ge n-4$. Let $A_0$ be the set of $a\in A$ such that
$\dd as\ne\dd b{s'}$ for all $s,s'\in S$, and $a\notin\Sph(s,\rho_1)\cup\Sph(s,\rho_2)$ for all $s\in S'$. By
Lemma~\ref{lem:ext} and the cap condition, $|A\setminus A_0|<j^2\sigma+2(n-4)\sigma$. It remains to show
$|A_0|\le e_A$.

\emph{Step 1: the $(n-2)$-set $T$.} For $a\in A_0$ let $T=T(a):=\{b,a\}\cup S'$, so $|T|=n-2$. Since
$\dd ab\notin R_b$, $a\notin\bigcup_s\Bis(b,s)$ and $\dd as\ne\dd b{s'}$, Lemma~\ref{lem:ext} shows that
$S\cup\{b,a\}$ is rainbow, hence so is $T$. Moreover $\rho_\nu\notin\Dist(T)$, where
$\Dist(T)=\Dist(S')\cup\{\dd bs,\dd as:s\in S'\}\cup\{\dd ab\}$. Indeed, $\rho_\nu\notin R_b\supseteq\Dist(S')\cup\{\dd bs\}$
by the definition of $\Bg_b$; $\rho_\nu\ne\dd ab$ since $a\in A$; and $\rho_\nu\ne\dd as$ for $s\in S'$ by the
definition of $A_0$.

\emph{Step 2: non-generic points.} Fix $a\in A_0$ and $\nu\in\{1,2\}$, and let $x\in C_\nu$. Then $x\notin T$:
$x\ne b$ since $\dd xb=\rho_\nu>0$, $x\ne a$ since $\dd ab\ne\rho_\nu$, and $x\notin S$ since $W\cap S=\emptyset$.
The label set of $x$ is $L_x=\{\rho_\nu,\dd xa\}\cup\{\dd xs:s\in S'\}$. If $x$ is not $T$-generic, then $x$ lies in
\[
\Sph\bigl(a,\Dist(T)\cup\{\rho_\nu\}\bigr)\ \cup\ \bigcup_{s\in S'}\Bis(a,s)\ \cup\ \Sm_\nu(a),\qquad
\Sm_\nu(a):=\bigcup_{s\in S'}\Sph\bigl(s,\{\rho_\nu,\dd ab\}\cup\{\dd b{s'},\dd a{s'}:s'\in S'\}\bigr).
\]
Indeed, the possible coincidences among $\Dist(T)\cup L_x$ are the following. $\Dist(T)$ itself is rainbow.
Equalities among the $\dd xs$, or between some $\dd xs$ and $\Dist(S')$, are impossible because $S\cup\{x\}$ is rainbow.
Equalities $\dd xs\in\{\rho_\nu,\dd ab,\dd b{s'},\dd a{s'}\}$ put $x$ in $\Sm_\nu(a)$. An equality $\dd xs=\dd xa$ puts $x$
in $\Bis(a,s)$. An equality $\dd xa\in\Dist(T)\cup\{\rho_\nu\}$ puts $x$ on a sphere about $a$. Finally
$\rho_\nu\notin\Dist(T)$ by Step~1. The set $\Sm_\nu(a)$ is a union of $(n-4)(2n-6)$ spheres about points of $S$,
so $|C_\nu\cap\Sm_\nu(a)|<2n^2\sigma$.

\emph{Step 3: counting cross pairs (``Lemma S2'').} Fix $a\in A_0$ and put $P:=C_1\cap Y_T$ and $Q':=C_2\cap Y_T$. These
are disjoint, since $\rho_1\ne\rho_2$. Let $E$ be the set of pairs $(x,w)\in P\times Q'$ with $\dd xa=\dd wa$, or
$\dd xa=\rho_2$, or $\dd wa=\rho_1$. We claim
\begin{equation}\label{eq:S2}
|(P\times Q')\setminus E|\le\bigl(4^{n-2}\mu+n(n-4)\sigma\bigr)(|P|+|Q'|).
\end{equation}
By Corollary~\ref{cor:G1}(b), at most $4^{n-2}\mu(|P|+|Q'|)$ pairs are label-disjoint. Let $(x,w)\notin E$ share a
label, where $L_x=\{\rho_1,\dd xa\}\cup\{\dd xs\}_{s\in S'}$ and $L_w=\{\rho_2,\dd wa\}\cup\{\dd ws\}_{s\in S'}$.
Since $\rho_1\ne\rho_2$ and $(x,w)\notin E$, the shared value involves some $s\in S'$, so one of the following
holds. Either $\dd xs\in\{\dd w{s'},\rho_2,\dd wa\}$: then for fixed $w$ the point $x$ lies on one of $(n-4)(n-2)$
spheres about points of $S'$, which gives at most $(n-4)(n-2)\sigma|Q'|$ pairs. Or $\dd ws\in\{\rho_1,\dd xa\}$:
this gives at most $2(n-4)\sigma|P|$ pairs. Since $(n-4)(n-2)\le n(n-4)$ and $2(n-4)\le n(n-4)$, \eqref{eq:S2}
follows.

\emph{Step 4: one of the two classes is covered.} Fix $a\in A_0$ and put $\theta:=Q^{1/2}\sigma$. We claim that there are
$\nu\in\{1,2\}$ and a real $t^\ast=t^\ast(a)$ such that, with $V(a):=\Dist(T)\cup\{\rho_1,\rho_2,t^\ast\}$,
\begin{equation}\label{eq:stepA-cover}
\Bigl|C_\nu\setminus\Bigl(\Sph(a,V(a))\cup\bigcup_{s\in S'}\Bis(a,s)\Bigr)\Bigr|\le\frac{|C_\nu|}{16(D+2)^{k'}} .
\end{equation}
We use $|C_\nu|\ge Q\sigma$ and the bounds $Q^{-1/2}=2^{-22}n^{-2}(D+2)^{-4k'}\le2^{-5}(D+2)^{-k'}$ and
$2n^2/Q\le2^{-5}(D+2)^{-k'}$.
\begin{itemize}[leftmargin=1.5em]
\item If $|P|<\theta$, Step~2 shows that the left side of~\eqref{eq:stepA-cover} for $\nu=1$ is at most
  $|P|+|C_1\cap\Sm_1(a)|<(Q^{1/2}+2n^2)\sigma\le(Q^{-1/2}+2n^2/Q)|C_1|\le|C_1|/(16(D+2)^{k'})$. The case $|Q'|<\theta$
  is symmetric, with $\nu=2$.
\item Otherwise $|P|,|Q'|\ge\theta$. Put $\eta:=(4^{n-2}\mu+n(n-4)\sigma)(1/|P|+1/|Q'|)$. Since
  $4^n\mu\le\smin\le\sigma$, we have $4^{n-2}\mu\le\sigma/16$, and therefore
  \[
  \eta\le\frac{2(\sigma/16+n^2\sigma)}{Q^{1/2}\sigma}\le 3n^2Q^{-1/2}\le2^{-20}(D+2)^{-4k'}\le\tfrac1{16}.
  \]
  Let $p_t:=|P\cap\Sph(a,t)|/|P|$ and $q_t:=|Q'\cap\Sph(a,t)|/|Q'|$. For independent uniform $x\in P$ and $w\in Q'$,
  \[
  \Pr[(x,w)\notin E]=\Pr[\dd xa\ne\rho_2,\ \dd wa\ne\rho_1,\ \dd xa\ne\dd wa]
  =(1-p_{\rho_2})(1-q_{\rho_1})-\sum_{t\notin\{\rho_1,\rho_2\}}p_tq_t,
  \]
  and by~\eqref{eq:S2} this is at most $\eta$. Lemma~\ref{lem:twodist} (with $\rho=\rho_1$, $\rho'=\rho_2$) gives
  $t^\ast$ with $p_{\rho_2}+p_{t^\ast}\ge1-\eta^{1/4}$ or $q_{\rho_1}+q_{t^\ast}\ge1-\eta^{1/4}$. In the first case all but
  $\eta^{1/4}|P|$ points of $P$ lie on $\Sph(a,\{\rho_2,t^\ast\})$. Together with Step~2, the left side
  of~\eqref{eq:stepA-cover} for $\nu=1$ is at most $\eta^{1/4}|C_1|+2n^2\sigma\le(2^{-5}(D+2)^{-k'}+2n^2/Q)|C_1|$,
  which is at most $|C_1|/(16(D+2)^{k'})$. The second case is symmetric, with $\nu=2$.
\end{itemize}

\emph{Step 5: GL, LR and OR.} Let $A_\nu$ be the set of $a\in A_0$ for which~\eqref{eq:stepA-cover} holds with this
$\nu$. Then $A_0=A_1\cup A_2$; say $|A_1|\ge|A_0|/2$ (the other case is symmetric). Note $A_1\cap C_1=\emptyset$.
List $V(a)$ as $v_1(a),\dots,v_{t+3}(a)$, where $t=\binom{n-2}2$ and entries are repeated if $|V(a)|<t+3$. On rows
$C_1$ and columns $A_1$ consider
\[
\mathcal B[x,a]:=\prod_{i=1}^{t+3}\bigl(\dd xa^2-v_i(a)^2\bigr)\prod_{s\in S'}\bigl(\dd xa^2-\dd xs^2\bigr).
\]
It has $t+n-1\le k'$ factors, so $\rk\mathcal B\le(D+2)^{k'}$ by Lemmas~\ref{lem:distrank} and~\ref{lem:HP}. By
\eqref{eq:stepA-cover}, each column has at most $|C_1|/(16(D+2)^{k'})$ nonzero entries. Lemma~\ref{lem:GL} gives
$C''\times A''$ with $|C''|\ge|C_1|/2$, $|A''|\ge|A_1|/2$ and $\mathcal B=0$ on it. So on $C''\times A''$ the matrix
$[\dd xa^2]$, of rank at most $D+2$, is covered by at most $k'\le\kst$ types: the column functions $v_i(a)^2$ (some of
them constants) and the row functions $\dd xs^2$, $s\in S'$. Theorem~\ref{thm:LR} gives $R\times C'\subseteq C''\times A''$
with $|R|\ge2^{-\Lst}|C''|$, $|C'|\ge2^{-\Lst}|A''|$ and a single type. Since $S'\cap C'=\emptyset$,
Lemma~\ref{lem:typeOR} gives $\min(|R|,|C'|)\le h$. But
$|R|\ge2^{-\Lst-1}Q\sigma\ge2^{-\Lst-1}\smin>h$, so $|C'|\le h$. Hence $|A''|\le2^{\Lst}h$, $|A_1|\le2^{\Lst+1}h$ and
$|A_0|\le2^{\Lst+2}h<e_A$.
\end{proof}

\subsection{Proof of Theorem~\ref{thm:Rplus}}
Claim$(n-2)$ was proved above. Let $j\in\{n-3,n-4\}$ and assume Claim$(j+1)$ with constant $\Gamma_{j+1}$. Let
$(S,W,\sigma)$ be a level-$j$ configuration with $\sigma\ge\smin$, and put $\tau:=Q\sigma\ge\smin$. For $b\in W$ define
\begin{gather*}
F_b:=\Delta_b\prod_{\lambda\in\Bg_b(\tau)}\bigl(\dd yb^2-\lambda^2\bigr)\ \text{ if }|\Bg_b(\tau)|\le1\ \text{(case B)},\\
F_b:=\Delta_b\,\bigl(\dd yb^2-\rho_1^2\bigr)\bigl(\dd yb^2-\rho_2^2\bigr)\ \text{ otherwise (case A)},
\end{gather*}
where in case A, $\rho_1\ne\rho_2$ are two elements of $\Bg_b(\tau)$. In both cases $F_b$ is a product of at most
$e_j=2+2j+\binom j2$ affine functions of the lifted coordinates, and $F_b(b)\ne0$.

\emph{Claim: for each $b\in W$, at most $\max(\Gamma_{j+1}Q\sigma,e_A)+(j^2+2n)\sigma$ points $x\in W\setminus\{b\}$
have $F_b(x)\ne0$.} In case B, the proof of the claim in Lemma~\ref{lem:kill} (with $K=1$ and $\tau=Q\sigma$) gives
fewer than $\Gamma_{j+1}Q\sigma+j^2\sigma$ such points. In case A, $F_b(x)\ne0$ implies $\dd xb\notin\{\rho_1,\rho_2\}$
and $\Delta_b(x)\ne0$, i.e.\ $x\in A$, and Lemma~\ref{lem:stepA} gives at most $e_A+(j^2+2(n-4))\sigma$ points.

As in the proof of Lemma~\ref{lem:kill}, Caro--Wei and Lemma~H give
$|W|\le P^+_j\bigl(2\max(\Gamma_{j+1}Q\sigma,e_A)+2(j^2+2n)\sigma+1\bigr)$. Since $e_A\le\sigma\le Q\Gamma_{j+1}\sigma$
and $1\le\sigma$ by~\eqref{eq:absorb}, we get $|W|\le P^+_j(2Q\Gamma_{j+1}+2(j^2+2n)+1)\sigma=\Gamma_j\sigma$. \qed

\begin{remark}
For $n=4$ the levels $j=2,1,0$ are all covered by Theorem~\ref{thm:Rplus}. In that case $S'=\emptyset$, the sets
$\Sm_\nu(a)$ are empty, and the cost $n(n-4)\sigma$ in~\eqref{eq:S2} vanishes. Claim$(0)$ then follows directly,
and Sections~\ref{sec:core} and~\ref{sec:slicing} are not needed.
\end{remark}

% ====================================================================
\section{Rainbow families and the Core Lemma}\label{sec:core}
% ====================================================================

Throughout Sections~\ref{sec:core} and~\ref{sec:slicing}, $n\ge5$, $0\le r\le n-5$, and $(S,W,\sigma)$ is a fixed level-$r$
configuration with $\sigma\ge\smin$. We assume that the conclusion of Lemma~\ref{lem:HR} holds for $(S,W,\sigma)$
with some constant $\Gamma'\ge1$; we refer to this as \textbf{HR$_r$}. For $c\in W$, $R_c=\Dist(S)\cup\{\dd cs:s\in S\}$
has at most $\binom r2+r$ elements. Recall from~\eqref{eq:absorb} that $\mu\le e\le\sigma$.

\subsection{Rainbow families}

\begin{definition}[rainbow family at level $r$]\label{def:family}
A \emph{rainbow family} $\mathcal F=(B,(C_k,\beta(k))_{k\in I})$ consists of a set $B=\{b_1,\dots,b_p\}\subseteq X$
with $p\ge1$, a finite index set $I$, and for each $k\in I$ a nonempty set $C_k\subseteq W$ and a $p$-tuple
$\beta(k)=(\beta_1(k),\dots,\beta_p(k))$ of positive reals, such that:
\begin{enumerate}[label=(F\arabic*)]
\item $S\cap B=\emptyset$ and $S\cup B$ is rainbow;
\item $\dd x{b_i}=\beta_i(k)$ for all $k\in I$, $x\in C_k$ and $i\le p$;
\item the $p|I|$ numbers $\beta_i(k)$ ($i\le p$, $k\in I$) are pairwise distinct, and none of them lies in
  $\Dist(S\cup B)$;
\item for all $k\in I$ and $x\in C_k$: $x\notin S\cup B$ and $S\cup B\cup\{x\}$ is rainbow;
\item $\dd xs\notin\beta(I):=\{\beta_i(k'):i\le p,\ k'\in I\}$ for all $x\in\bigcup_kC_k$ and $s\in S$.
\end{enumerate}
We write $M(\mathcal F):=\min_k|C_k|$, and also use $\beta(k)$ for the \emph{set} of its $p$ values.
\end{definition}

The classes $C_k$ are pairwise disjoint, since they lie on spheres about $b_1$ with distinct radii
$\beta_1(k)$, and they are disjoint from $S\cup B$. Passing to a subset $I'\subseteq I$, and to nonempty subsets
$C'_k\subseteq C_k$, preserves (F1)--(F5).

\subsection{Constants}
Let $p:=n-4-r$ and let $t,w_1,w_2,k_3,k_4,k_5,k_c$ be as in Definition~\ref{def:typecounts}. Put
\begin{gather*}
\theta_{\mathrm{rel}}(w):=\frac{1}{32\,\Gamma'(D+2)^{k_4(w)}},\qquad
\eps:=\frac{\theta_{\mathrm{rel}}(w_2)}{96\,(D+2)^{k_3(w_2)}},\\
\hat M_{\mathrm{core}}:=2^{20}4^nn^2\Gamma'^2(D+2)^{2k_c},\qquad M_{\mathrm{core}}:=\hat M_{\mathrm{core}}\,\sigma .
\end{gather*}

\begin{theorem}[Core Lemma]\label{thm:core}
Assume HR$_r$ with constant $\Gamma'$. There is no rainbow family at level $r$ with $p=|B|=n-4-r$, $|I|\ge12$ and
$M(\mathcal F)\ge M_{\mathrm{core}}$.
\end{theorem}

The rest of this section proves Theorem~\ref{thm:core}. Suppose $\mathcal F$ is such a family; discarding indices
we may assume $|I|=12$. Put $M:=M(\mathcal F)\ge M_{\mathrm{core}}$ and
\[
\delta_1:=\frac{2\cdot4^{n-2}\mu}{\eps M}+\frac{3n^2\sigma}{M}.
\]

\begin{lemma}[parameter inequalities]\label{lem:params-core}
\begin{enumerate}[label=(P\arabic*)]
\item $\delta_1\le\theta_{\mathrm{rel}}(w_1)/(16(D+2)^{k_3(w_1)})=1/(2^9\Gamma'(D+2)^{k_c})$;
\item $2n^2\sigma\le\eps M$;
\item $6\eps=\theta_{\mathrm{rel}}(w_2)/(16(D+2)^{k_3(w_2)})$;
\item $M\ge32(D+2)^{k_4(w_1)}\Gamma'(r+1)\sigma$;
\item $M\ge80n^2(D+2)^{k_5}\sigma$;
\item $M\ge\sigma\ge e>2^{\Lst+1}h$.
\end{enumerate}
\end{lemma}
\begin{proof}
(P3) is the definition of $\eps$, and (P6) follows from~\eqref{eq:absorb}. We have
$1/\eps=96\cdot32\,\Gamma'(D+2)^{k_3(w_2)+k_4(w_2)}\le3\cdot2^{10}\Gamma'(D+2)^{k_c}$. For (P1), the first term of
$\delta_1$ is at most $6\cdot2^{10}\,4^{n-2}\Gamma'(D+2)^{k_c}\sigma/M$, using $\mu\le\sigma$. This is at most
$1/(2^{10}\Gamma'(D+2)^{k_c})$ because $M\ge M_{\mathrm{core}}\ge6\cdot2^{20}4^{n-2}\Gamma'^2(D+2)^{2k_c}\sigma$. The
second term is at most $1/(2^{10}\Gamma'(D+2)^{k_c})$ because $M\ge3\cdot2^{10}n^2\Gamma'(D+2)^{k_c}\sigma$.
(P2) holds because $M\ge6\cdot2^{10}n^2\Gamma'(D+2)^{k_c}\sigma\ge2n^2\sigma/\eps$. (P4) and (P5) hold because
$k_4(w_1),k_5\le k_c$, $r+1\le n$ and $\Gamma'\ge1$.
\end{proof}

\subsection{Notation for two auxiliary points}
Fix $l\ne m$ in $I$. For $a_1\in C_l$ and $a_2\in C_m$ put $T=T(a_1,a_2):=S\cup B\cup\{a_1,a_2\}$. These are $n-2$
distinct points, since $r+p=n-4$, $a_1\ne a_2$, and $a_1,a_2\notin S\cup B$ by (F4). For $x\notin T$ let
$L_x:=\{\dd xy:y\in T\}$. For $x\in C_k$ we have $L_x=\{\dd xs\}_{s\in S}\cup\beta(k)\cup\{\dd x{a_1},\dd x{a_2}\}$ by (F2).
Let $Y_T$ be the set of $T$-generic points, and put
\[
\Hyp(a_1,a_2):=\bigcup_{s\in S,\,i=1,2}\Bis(s,a_i),\qquad
\Sm(a_1,a_2):=\bigcup_{s\in S}\Sph\bigl(s,\{\dd{a_i}{s'}:i=1,2,\ s'\in S\}\cup\{\dd{a_1}{a_2}\}\bigr).
\]
Both are empty if $r=0$. $\Sm$ is a union of $r(2r+1)$ spheres about points of $S$, so it meets $W$ in fewer than
$r(2r+1)\sigma\le2n^2\sigma$ points. By (F2),
\begin{equation}\label{eq:DT}
\Dist(T)=\Dist(S\cup B)\cup\{\dd{a_i}s:i=1,2,\ s\in S\}\cup\beta(l)\cup\beta(m)\cup\{\dd{a_1}{a_2}\},\qquad|\Dist(T)|\le t .
\end{equation}

\begin{definition}[validity]\label{def:valid}
A pair $(k,k')$ is \emph{valid for $(a_1,a_2)$} if $k\ne k'$ are in $I\setminus\{l,m\}$, $T(a_1,a_2)$ is rainbow,
and $(\beta(k)\cup\beta(k'))\cap\Dist(T)=\emptyset$.
\end{definition}

Validity is used only in Steps~0--2, to obtain the inclusions there. The properties $\Good_k$ and $\Cov_k$ defined
below are properties of \emph{arbitrary} pairs $(a_1,a_2)\in C_l\times C_m$. Steps~3--5 count them over all of
$C_l\times C_m$, and validity enters Step~5 only through the set $V_{lm}$ of pairs on which it holds
(referee item~F).

\subsection{Step 0: non-generic loci}

\begin{lemma}[Step 0]\label{lem:step0}
Let $T$ be rainbow, $k\in I\setminus\{l,m\}$ and $\beta(k)\cap\Dist(T)=\emptyset$. Every $x\in C_k\setminus Y_T$ lies in
\[
\Sph\bigl(a_1,\Dist(T)\cup\beta(k)\bigr)\cup\Sph\bigl(a_2,\Dist(T)\cup\beta(k)\bigr)\cup\Bis(a_1,a_2)\cup\Hyp(a_1,a_2)\cup\Sm(a_1,a_2).
\]
\end{lemma}
\begin{proof}
First, $x\notin T$, because $C_k$ is disjoint from $S\cup B$, $C_l$ and $C_m$. We list the possible coincidences
among $\Dist(T)\cup L_x$. $\Dist(T)$ is rainbow. Equalities among the $\dd xs$ and $\beta(k)$, or between these
and $\Dist(S\cup B)$, are excluded by (F4) and (F3). The case $\dd xs\in\beta(l)\cup\beta(m)$ is excluded by (F5).
Equalities $\dd xs\in\{\dd{a_i}{s'},\dd{a_1}{a_2}\}$ give $\Sm$. An equality $\dd xs=\dd x{a_i}$ gives $\Hyp$. An
equality $\dd x{a_i}\in\Dist(T)\cup\beta(k)$ gives a sphere about $a_i$, and $\dd x{a_1}=\dd x{a_2}$ gives
$\Bis(a_1,a_2)$. Finally $\beta(k)\cap\Dist(T)=\emptyset$ by hypothesis.
\end{proof}

\subsection{Step 1: a trichotomy}
For $k\in I\setminus\{l,m\}$ put
\[
\Lambda_k(a_1,a_2):=\Sph\bigl(a_1,\Dist(T)\cup\beta(k)\bigr)\cup\Sph\bigl(a_2,\Dist(T)\cup\beta(k)\bigr)\cup\Bis(a_1,a_2)\cup\Hyp(a_1,a_2).
\]

\begin{definition}[$\Good$]\label{def:good}
$\Good_k(a_1,a_2)$ holds if there are a set $\Omega$ of at most $w_1$ positive reals and a set $\Ex\subseteq C_k$
with $|\Ex|\le\delta_1|C_k|$ such that
$C_k\subseteq\Sph(a_1,\Omega)\cup\Sph(a_2,\Omega)\cup\Bis(a_1,a_2)\cup\Hyp(a_1,a_2)\cup\Ex$.
\end{definition}

\begin{lemma}[Step 1]\label{lem:step1}
Let $(k,k')$ be valid for $(a_1,a_2)$. Then $\Good_k(a_1,a_2)$ holds, or
\[
\bigl|C_{k'}\setminus\bigl(\Lambda_{k'}(a_1,a_2)\cup\Sph(a_1,\beta(k))\cup\Sph(a_2,\beta(k))\bigr)\bigr|<2\eps|C_{k'}|.
\tag*{(ii)$_{k,k'}$}
\]
\end{lemma}
\begin{proof}
Put $P:=C_k\cap Y_T$, $Q:=C_{k'}\cap Y_T$ and $Q':=\{w\in Q:L_w\cap\beta(k)=\emptyset\}$.

\emph{Case $|Q'|\ge\eps|C_{k'}|$.} $P$ and $Q'$ are disjoint subsets of $Y_T$. By Corollary~\ref{cor:G1}(b), at most
$4^{n-2}\mu(|C_k|+|C_{k'}|)$ pairs $(w,x)\in Q'\times P$ are label-disjoint. So some $w_0\in Q'$ is label-disjoint
from at most $\mathrm{ex}:=4^{n-2}\mu(|C_k|+|C_{k'}|)/(\eps|C_{k'}|)$ points of $P$. Let $x\in P$ be any other point,
so that $L_x\cap L_{w_0}\ne\emptyset$. Now $\beta(k)\cap L_{w_0}=\emptyset$, because $\beta(k)\cap\beta(k')=\emptyset$ (F3),
$\dd{w_0}s\notin\beta(k)$ (F5), and $w_0\in Q'$. So the common value is $\dd xs$ for some $s$, or $\dd x{a_i}$
for some $i$, and it lies in $L_{w_0}$. If it is $\dd xs$, then $\dd xs\notin\beta(k')$ by (F5), so
$\dd xs\in\{\dd{w_0}{s'}\}_{s'}\cup\{\dd{w_0}{a_1},\dd{w_0}{a_2}\}$, and $x$ lies on one of $r(r+2)$ spheres about
points of $S$. If it is $\dd x{a_i}$, then $x\in\Sph(a_i,L_{w_0})$. Combining this with Lemma~\ref{lem:step0}, $C_k$ is
covered as in Definition~\ref{def:good} with $\Omega:=L_{w_0}\cup\Dist(T)\cup\beta(k)$, where
$|\Omega|\le(n-2)+t+p=w_1$, and with $\Ex$ consisting of the at most $\mathrm{ex}$ label-disjoint partners, the
fewer than $r(r+2)\sigma$ points on the spheres just described, and the fewer than $r(2r+1)\sigma$ points of
$C_k\cap\Sm$. Since $|C_k|,|C_{k'}|\ge M$,
\[
|\Ex|\le\Bigl(\frac{4^{n-2}\mu}{\eps}\Bigl(\frac1{|C_{k'}|}+\frac1{|C_k|}\Bigr)+\frac{3n^2\sigma}{M}\Bigr)|C_k|\le\delta_1|C_k|,
\]
so $\Good_k(a_1,a_2)$ holds.

\emph{Case $|Q'|<\eps|C_{k'}|$.} By Lemma~\ref{lem:step0} applied to $k'$ (valid since $\beta(k')\cap\Dist(T)=\emptyset$),
$C_{k'}\setminus Q\subseteq\Lambda_{k'}\cup\Sm$. If $w\in Q\setminus Q'$, then $L_w\cap\beta(k)\ne\emptyset$. Since
$\dd ws\notin\beta(k)$ (F5) and $\beta(k')\cap\beta(k)=\emptyset$ (F3), this means $\dd w{a_i}\in\beta(k)$ for some
$i$, i.e.\ $w\in\Sph(a_1,\beta(k))\cup\Sph(a_2,\beta(k))$. Hence the set in (ii)$_{k,k'}$ is contained in
$Q'\cup(C_{k'}\cap\Sm)$, which has fewer than $\eps|C_{k'}|+2n^2\sigma\le2\eps|C_{k'}|$ points, by (P2).
\end{proof}

\subsection{Step 2: summing over $k$}

\begin{lemma}[Step 2]\label{lem:step2}
Fix $a_1,a_2$ and $k'\in I\setminus\{l,m\}$. Let $J\subseteq I\setminus\{l,m,k'\}$ with $|J|\ge3$ be such that, for every
$k\in J$, $(k,k')$ is valid and (ii)$_{k,k'}$ holds. Then
\[
\Cov_{k'}(a_1,a_2):\qquad |C_{k'}\setminus\Lambda_{k'}(a_1,a_2)|<6\eps|C_{k'}| .
\]
\end{lemma}
\begin{proof}
Let $c:=|C_{k'}\setminus\Lambda_{k'}|$ and assume $c\ge2\eps|C_{k'}|$. For each $k\in J$, by (ii)$_{k,k'}$, more than
$c-2\eps|C_{k'}|$ points of $C_{k'}\setminus\Lambda_{k'}$ lie in $\Sph(a_1,\beta(k))\cup\Sph(a_2,\beta(k))$. By (F3) the
radius sets $\beta(k)$, $k\in J$, are pairwise disjoint, so the sets $\Sph(a_1,\beta(k))$, $k\in J$, are pairwise
disjoint, and likewise for $a_2$. Summing over $J$ gives $|J|(c-2\eps|C_{k'}|)<2c$, i.e.\
$c(|J|-2)<2|J|\eps|C_{k'}|$, so $c<6\eps|C_{k'}|$ as $|J|\ge3$.
\end{proof}

Note that $\Cov_{k'}(a_1,a_2)$ has the shape of $\Good_{k'}(a_1,a_2)$ with $\Omega:=\Dist(T)\cup\beta(k')$, so that
$|\Omega|\le t+p=w_2$, and $\Ex:=C_{k'}\setminus\Lambda_{k'}$, so that $|\Ex|<6\eps|C_{k'}|$. Here $\Omega$ depends on $(a_1,a_2)$
through $\Dist(T)$.

\subsection{Step 3: the two-centre lemma}

\begin{definition}[$w$-few]\label{def:wfew}
Let $j\in I\setminus\{l,m\}$, $w\in\{w_1,w_2\}$ and $a_2\in C_m$. Put $\theta_w:=\max(\theta_{\mathrm{rel}}(w)|C_j|,\sigma)$ and
$\theta'_w:=\theta_w+r\sigma$. We say that \emph{$a_2$ is $w$-few on $C_j$} if there is a set $\Xi$ of at most
$w+r+\binom r2$ positive reals with
\[
\Bigl|C_j\setminus\Bigl(\Sph(a_2,\Xi)\cup\bigcup_{s\in S}\Bis(a_2,s)\Bigr)\Bigr|\le\Gamma'\theta'_w .
\]
This property depends only on $(a_2,j,w)$.
\end{definition}

\begin{lemma}[Step 3, two-centre lemma]\label{lem:step3}
Let $j\in I\setminus\{l,m\}$, $a_2\in C_m$, $w\in\{w_1,w_2\}$ and $0<\delta\le\theta_{\mathrm{rel}}(w)/(16(D+2)^{k_3(w)})$. Let
$Z\subseteq C_l$, and suppose that for each $a_1\in Z$ there are a set $\Omega(a_1)$ of at most $w$ positive reals and a
set $\Ex(a_1)\subseteq C_j$ with $|\Ex(a_1)|\le\delta|C_j|$ such that
\[
C_j\subseteq\Sph(a_1,\Omega(a_1))\cup\Sph(a_2,\Omega(a_1))\cup\Bis(a_1,a_2)\cup\Hyp(a_1,a_2)\cup\Ex(a_1).
\]
If $a_2$ is not $w$-few on $C_j$, then $|Z|\le(w+1)e$.
\end{lemma}
\begin{proof}
For a real $\lambda>0$ let $R_\lambda:=(C_j\cap\Sph(a_2,\lambda))\setminus\bigcup_{s\in S}\Sph(s,\lambda)$ and
$Z_\lambda:=\{a_1\in Z:\lambda\notin\Omega(a_1)\}$. Let $a_1\in Z_\lambda$ and $x\in R_\lambda\setminus\Ex(a_1)$. Then
$x\notin\Sph(a_2,\Omega(a_1))$, since $\dd x{a_2}=\lambda\notin\Omega(a_1)$. Also $x\notin\Bis(s,a_2)$, since otherwise
$\dd xs=\lambda$. If $x\in\Bis(a_1,a_2)$ then $\dd x{a_1}=\lambda$, and if $x\in\Bis(s,a_1)$ then $\dd x{a_1}=\dd xs$.
Hence
\begin{equation}\label{eq:31}
\dd x{a_1}^2\in\{v^2:v\in\Omega(a_1)\}\cup\{\lambda^2\}\cup\{\dd xs^2:s\in S\}\qquad(x\in R_\lambda\setminus\Ex(a_1)).
\end{equation}
\emph{Claim: if $|R_\lambda|\ge\theta_w$, then $|Z_\lambda|\le e$.} Take rows $R_\lambda$ and columns $Z_\lambda$; they are
disjoint, as $C_j\cap C_l=\emptyset$. List $\Omega(a_1)$ as $v_1(a_1),\dots,v_w(a_1)$, repeating entries if necessary, and put
\[
\mathcal B[x,a_1]:=\prod_{i\le w}\bigl(\dd x{a_1}^2-v_i(a_1)^2\bigr)\cdot\bigl(\dd x{a_1}^2-\lambda^2\bigr)\cdot\prod_{s\in S}\bigl(\dd x{a_1}^2-\dd xs^2\bigr).
\]
This is a product of $k_3(w)=w+1+r$ factors, so $\rk\mathcal B\le(D+2)^{k_3(w)}$. By~\eqref{eq:31}, column $a_1$ has nonzero
entries only in $\Ex(a_1)$, hence at most $\delta|C_j|\le\theta_{\mathrm{rel}}(w)|C_j|/(16(D+2)^{k_3(w)})\le|R_\lambda|/(16(D+2)^{k_3(w)})$
of them. Lemma~\ref{lem:GL} gives $R''\times Z''$ with $|R''|\ge|R_\lambda|/2$, $|Z''|\ge|Z_\lambda|/2$ and $\mathcal B=0$ there.
On $R''\times Z''$ the matrix $[\dd x{a_1}^2]$ is covered by $k_3(w)\le\kst$ types: the column functions $v_i(a_1)^2$, the
constant $\lambda^2$, and the row functions $\dd xs^2$. Theorem~\ref{thm:LR} and Lemma~\ref{lem:typeOR} (note
$S\cap Z''=\emptyset$) give $R'''\times Z'''$ with $\min(|R'''|,|Z'''|)\le h$, where
$|R'''|\ge2^{-\Lst}|R_\lambda|/2\ge2^{-\Lst-1}\sigma>h$ by (P6). So $|Z'''|\le h$, hence $|Z''|\le2^{\Lst}h$ and
$|Z_\lambda|\le2^{\Lst+1}h\le e$.

Suppose first that there are $w+1$ distinct radii $\lambda_1,\dots,\lambda_{w+1}$ with $|R_{\lambda_i}|\ge\theta_w$. Since
$|\Omega(a_1)|\le w$, every $a_1\in Z$ lies in some $Z_{\lambda_i}$, so $|Z|\le(w+1)e$. Otherwise let $\Xi_h$ be the set of
the at most $w$ radii with $|R_\lambda|\ge\theta_w$, and put $\Xi:=\Xi_h\cup R_{a_2}$, so that $|\Xi|\le w+r+\binom r2$. Let
$x\in C_j$ with $x\notin\bigcup_s\Bis(a_2,s)$ and $\lambda:=\dd x{a_2}\notin\Xi$. Then $|R_\lambda|<\theta_w$, so
$|C_j\cap\Sph(a_2,\lambda)|<\theta_w+r\sigma=\theta'_w$, because the spheres $\Sph(s,\lambda)$, $s\in S$, meet $W$ in fewer than
$r\sigma$ points in total. Apply HR$_r$ with $c:=a_2\in W$, $V:=C_j\subseteq W\setminus\{a_2\}$ and threshold
$\theta'_w\ge\sigma$. Since $\lambda\notin R_{a_2}$, every such $x$ is counted there, so there are at most $\Gamma'\theta'_w$ of
them. Hence $a_2$ is $w$-few on $C_j$.
\end{proof}

So Lemma~\ref{lem:step3} says: \emph{if $a_2$ is not $w$-few on $C_j$, then $|Z|\le(w+1)e$} (referee item~X4).

\subsection{Step 4: few $w$-few centres}

\begin{lemma}[Step 4]\label{lem:step4}
For $j\in I\setminus\{l,m\}$ and $w\in\{w_1,w_2\}$, the set $F$ of points $a_2\in C_m$ that are $w$-few on $C_j$ has $|F|\le e$.
\end{lemma}
\begin{proof}
Take rows $C_j$ and columns $F$ (disjoint). For $a_2\in F$ let $\Xi(a_2)$ be as in Definition~\ref{def:wfew}, padded to
exactly $w+r+\binom r2$ entries, and put
$\mathcal B[x,a_2]:=\prod_{v\in\Xi(a_2)}(\dd x{a_2}^2-v^2)\prod_{s\in S}(\dd x{a_2}^2-\dd xs^2)$. This is a product of
$k_4(w)$ factors, so $\rk\mathcal B\le(D+2)^{k_4(w)}$. Column $a_2$ has at most $\Gamma'\theta'_w$ nonzero entries, and
\[
\Gamma'\theta'_w\le\Gamma'\theta_{\mathrm{rel}}(w)|C_j|+\Gamma'(r+1)\sigma\le\frac{|C_j|}{32(D+2)^{k_4(w)}}+\frac{|C_j|}{32(D+2)^{k_4(w)}},
\]
by the definition of $\theta_{\mathrm{rel}}$ and (P4) (as $k_4(w)\le k_4(w_1)$). Lemma~\ref{lem:GL} gives a zero rectangle
$C''\times F''$ with halves. There $[\dd x{a_2}^2]$ is covered by $k_4(w)\le\kst$ types: the column functions $a_2\mapsto v(a_2)^2$
given by the (padded) entries $v$ of $\Xi(a_2)$, and the row functions $\dd xs^2$. Theorem~\ref{thm:LR} and Lemma~\ref{lem:typeOR} give a rectangle with
$\min\le h$, whose row side is at least $2^{-\Lst}|C_j|/2>h$ by (P6). So $|F''|\le2^{\Lst}h$ and $|F|\le e$.
\end{proof}

\begin{corollary}\label{cor:step34}
Let $j\in I\setminus\{l,m\}$, $w\in\{w_1,w_2\}$ and $\delta$ be as in Lemma~\ref{lem:step3}, and suppose that for every
$a_2\in C_m$ the set $Z(a_2)\subseteq C_l$ satisfies the hypothesis of Lemma~\ref{lem:step3}. Then
$\sum_{a_2\in C_m}|Z(a_2)|\le e|C_l|+(w+1)e|C_m|$.
\end{corollary}
\begin{proof}
The at most $e$ points $a_2$ that are $w$-few on $C_j$ (Lemma~\ref{lem:step4}) contribute at most $|C_l|$ each. Every other
$a_2$ contributes at most $(w+1)e$, by Lemma~\ref{lem:step3}.
\end{proof}

We apply Corollary~\ref{cor:step34} in two ways, for each $j\in I\setminus\{l,m\}$:
\begin{enumerate}[label=(\alph*)]
\item $Z(a_2):=\{a_1\in C_l:\Good_j(a_1,a_2)\}$, with $w=w_1$ and $\delta=\delta_1$, which is admissible by (P1). This gives
  $|\{(a_1,a_2):\Good_j(a_1,a_2)\}|\le e|C_l|+(w_1+1)e|C_m|$.
\item $Z(a_2):=\{a_1\in C_l:\Cov_j(a_1,a_2)\}$, with $w=w_2$ and $\delta=6\eps$, which is admissible by (P3); see the
  remark after Lemma~\ref{lem:step2}. This gives $|\{(a_1,a_2):\Cov_j(a_1,a_2)\}|\le e|C_l|+(w_2+1)e|C_m|$.
\end{enumerate}

\subsection{Step 5: counting, and the contradiction}
Let $\Pi$ be the set of pairs $(a_1,a_2)\in C_l\times C_m$ with $\dd{a_1}s=\dd{a_2}{s'}$ for some $s,s'\in S$. For fixed $a_2$,
fewer than $r^2\sigma$ points $a_1$ qualify (cap), so $|\Pi|<r^2\sigma|C_m|$; and $\Pi=\emptyset$ if $r=0$. Put
\[
\Ty(a_1,a_2):=\Dist(S\cup B)\cup\beta(l)\cup\beta(m)\cup\{\dd{a_1}s:s\in S\}\cup\{\dd{a_2}s:s\in S\},
\]
a list of at most $\binom{r+p}2+2p+2r=k_5$ numbers, and $V_{lm}:=\{(a_1,a_2)\notin\Pi:\dd{a_1}{a_2}\notin\Ty(a_1,a_2)\}$.

\begin{lemma}\label{lem:51}
$T(a_1,a_2)$ is rainbow if and only if $(a_1,a_2)\notin\Pi$ and $\dd{a_1}{a_2}\notin\Ty(a_1,a_2)$.
\end{lemma}
\begin{proof}
The distances of $T$ are $\Dist(S\cup B)$, $\dd{a_1}s$, $\beta(l)$, $\dd{a_2}s$, $\beta(m)$ and $\dd{a_1}{a_2}$. The sets
$S\cup B\cup\{a_i\}$ are rainbow (F4), $\beta(l)\cap\beta(m)=\emptyset$ (F3), and $\dd{a_1}s\notin\beta(m)$ and
$\dd{a_2}s\notin\beta(l)$ (F5). So the only possible coincidences are $\dd{a_1}s=\dd{a_2}{s'}$ (that is, $\Pi$) and
$\dd{a_1}{a_2}\in\Ty$.
\end{proof}

Let $(a_1,a_2)\in V_{lm}$. Then $T$ is rainbow. By (F3) at most one index $k^\bullet\in I$ has $\dd{a_1}{a_2}\in\beta(k^\bullet)$,
and $k^\bullet\notin\{l,m\}$ since $\beta(l)\cup\beta(m)\subseteq\Ty$. Put $I_0:=I\setminus\{l,m,k^\bullet\}$, so $|I_0|\ge9$. For
$k\in I_0$ we have $\beta(k)\cap\Dist(T)=\emptyset$ by~\eqref{eq:DT}, (F3), (F5) and $k\ne k^\bullet$. Hence every
$(k,k')$ with $k\ne k'$ in $I_0$ is valid. Let $g:=|\{k\in I_0:\Good_k(a_1,a_2)\}|$. If $g\le|I_0|-4$, then for each
$k'\in I_0$ the set $J:=\{k\in I_0\setminus\{k'\}:\text{not }\Good_k(a_1,a_2)\}$ has $|J|\ge3$, and (ii)$_{k,k'}$ holds for
$k\in J$ by Lemma~\ref{lem:step1}; so $\Cov_{k'}(a_1,a_2)$ holds by Lemma~\ref{lem:step2}, for every $k'\in I_0$. In either case
\[
|\{k\in I\setminus\{l,m\}:\Good_k(a_1,a_2)\}|+|\{k\in I\setminus\{l,m\}:\Cov_k(a_1,a_2)\}|\ge6\qquad((a_1,a_2)\in V_{lm}).
\]
Summing over $(a_1,a_2)\in C_l\times C_m$ and applying (a) and (b) to each of the $10$ indices $j\in I\setminus\{l,m\}$,
\[
6|V_{lm}|\le10\bigl[2e|C_l|+(w_1+w_2+2)e|C_m|\bigr]\le10(w_1+w_2+2)\,e\,(|C_l|+|C_m|).
\]
So the density of $V_{lm}$ in $C_l\times C_m$ is at most $\frac{10}6(w_1+w_2+2)e(\frac1{|C_m|}+\frac1{|C_l|})\le4(w_1+w_2+2)e/M$
(referee item~B: the factor $10/6$ is the ratio of the number of indices summed, $10$, to the guaranteed count, $6$).
The density of $\Pi$ is less than $r^2\sigma/|C_l|\le r^2\sigma/M$. Since
$w_1+w_2+2=n^2-2n-2-2r<n^2$ and $e\le\sigma$, the set $V_{lm}\cup\Pi$ has density at most $5n^2\sigma/M\le1/(16(D+2)^{k_5})$, by (P5).

Now consider $\mathcal B_0[a_1,a_2]:=\prod_{v\in\Ty(a_1,a_2)}(\dd{a_1}{a_2}^2-v^2)$ on $C_l\times C_m$, with $\Ty$ padded to $k_5$
entries. It has rank at most $(D+2)^{k_5}$, and it vanishes outside $V_{lm}\cup\Pi$. Lemma~\ref{lem:GL} (average form) gives
$C_l''\times C_m''$ with halves and $\mathcal B_0=0$. There $[\dd{a_1}{a_2}^2]$ is covered by $k_5\le\kst$ types: the constants
in $\Dist(S\cup B)\cup\beta(l)\cup\beta(m)$, the row functions $\dd{a_1}s^2$ and the column functions $\dd{a_2}s^2$. By
Theorem~\ref{thm:LR} there is a single-type rectangle $R\times C$ with $|R|\ge2^{-\Lst}|C_l|/2>h$ and
$|C|\ge2^{-\Lst}|C_m|/2>h$, using (P6). Lemma~\ref{lem:typeOR} (rows $a_1\in R$, columns $a_2\in C$, $S\cap C=\emptyset$) gives
$\min(|R|,|C|)\le h$, a contradiction. This proves Theorem~\ref{thm:core}. \qed

\begin{remark}
For $n=5$ and $r=0$ (so $p=1$) the sets $\Hyp$, $\Sm$ and $\Pi$ are empty, $R_{a_2}=\emptyset$, and only column functions
and constants occur. In this case the proof is the argument for $f_d(5)$ written in the project's round~13 and
refereed separately (``referee14''), up to the values of the constants. For general $r$ the proof adds
row types $\dd xs^2$, coming from the bisectors $\Bis(s,a_i)$, and the small classes of points of $S$, whose
costs are absorbed in $\delta_1$, in $\eps$ and in Step~5.
\end{remark}

% ====================================================================
\section{The Slicing Lemma and the class-count theorem E1$_r$}\label{sec:slicing}
% ====================================================================

The setting is that of Section~\ref{sec:core}: $n\ge5$, $0\le r\le n-5$, $(S,W,\sigma)$ is a level-$r$ configuration with
$\sigma\ge\smin$, and HR$_r$ holds for it with constant $\Gamma'\ge1$.

\subsection{The Slicing Lemma}

\begin{lemma}[Slicing Lemma]\label{lem:slicing}
Let $0\le r\le n-6$, let $K\ge1$, and let $\mathcal F=(B,(C_k,\beta(k))_{k\in I})$ be a rainbow family at level $r$
with $1\le p=|B|\le n-5-r$ and $|I|\ge6K$. Put
\[
k_{\mathrm{sl}}:=3K-1+\tbinom{r+p}2+2p+2r,\qquad\theta_{\mathrm{sl}}:=\frac1{16\,\Gamma'(D+2)^{k_{\mathrm{sl}}}},\qquad
M_{\mathrm{sl}}:=\frac{2n(n+K)\sigma}{\theta_{\mathrm{sl}}},
\]
and assume $k_{\mathrm{sl}}\le\kst$. If $M(\mathcal F)\ge M_{\mathrm{sl}}$, then there are a point $c\in\bigcup_kC_k$ and a rainbow family
$\mathcal F'=(B\cup\{c\},(C'_k,\beta'(k))_{k\in I'})$ at level $r$ with $I'\subseteq I$, $|I'|=K$, $C'_k\subseteq C_k$ and
$M(\mathcal F')\ge\theta_{\mathrm{sl}}M(\mathcal F)/2$.
\end{lemma}

\begin{proof}
Discard indices so that $|I|=6K$, and put $M:=M(\mathcal F)$. Fix $\pi\in I$. For $c\in C_\pi$ the set $S\cup B\cup\{c\}$ is
rainbow by (F4), and since $\dd c{b_i}=\beta_i(\pi)$,
\[
\Phi(c):=\Dist(S\cup B\cup\{c\})=\Dist(S\cup B)\cup\beta(\pi)\cup\{\dd cs:s\in S\},\qquad|\Phi(c)|\le\tbinom{r+p}2+p+r .
\]
Note that $R_c=\Dist(S)\cup\{\dd cs\}\subseteq\Phi(c)$. Let $G_c$ be the bipartite graph with left vertex set $I\setminus\{\pi\}$ and
right vertex set $\R_{>0}$, in which $k\sim\lambda$ if $\lambda\notin\Phi(c)\cup\beta(k)$ and
$|C_k\cap\Sph(c,\lambda)|\ge\theta_{\mathrm{sl}}|C_k|$. Each $k$ has at most $1/\theta_{\mathrm{sl}}$ neighbours, so only finitely many
right vertices are non-isolated, and we regard $G_c$ as a finite graph.

\emph{(a) Few bad centres.} Let $Z_{\mathrm{bad}}$ be the set of $c\in C_\pi$ for which $G_c$ has no matching of size $3K$.
We show $|Z_{\mathrm{bad}}|\le2e$. For $c\in Z_{\mathrm{bad}}$, Theorem~\ref{thm:Konig} gives a vertex cover
$\mathcal A(c)\cup\Lambda_0(c)$, with $\mathcal A(c)\subseteq I\setminus\{\pi\}$, $\Lambda_0(c)\subseteq\R_{>0}$ and
$|\mathcal A(c)|+|\Lambda_0(c)|\le3K-1$. Call $c$ \emph{$k$-few} for every $k\in I\setminus(\{\pi\}\cup\mathcal A(c))$; then every
edge at $k$ ends in $\Lambda_0(c)$.

Let $c$ be $k$-few, and let $x\in C_k$ satisfy $x\notin\bigcup_s\Bis(c,s)$ and
$\lambda:=\dd xc\notin\Lambda_0(c)\cup\Phi(c)\cup\beta(k)$. Then $(k,\lambda)$ is not an edge, since $\lambda\notin\Lambda_0(c)$.
As $\lambda\notin\Phi(c)\cup\beta(k)$, this means $|C_k\cap\Sph(c,\lambda)|<\theta_{\mathrm{sl}}|C_k|$. Also $\lambda\notin R_c$. Now apply
HR$_r$ with $c\in W$, $V:=C_k\subseteq W\setminus\{c\}$ and threshold $\theta:=\theta_{\mathrm{sl}}|C_k|$; this is admissible since
$\theta\ge\theta_{\mathrm{sl}}M_{\mathrm{sl}}=2n(n+K)\sigma\ge\sigma$. So the number of such $x$ is at most
$\Gamma'\theta_{\mathrm{sl}}|C_k|=|C_k|/(16(D+2)^{k_{\mathrm{sl}}})$. Hence the matrix
\[
\mathcal B_k[x,c]:=\prod_{v\in\Lambda_0(c)\cup\Phi(c)\cup\beta(k)}\bigl(\dd xc^2-v^2\bigr)\prod_{s\in S}\bigl(\dd xc^2-\dd xs^2\bigr)
\]
(with $\Lambda_0(c)$ padded to $3K-1$ entries, so that there are at most $k_{\mathrm{sl}}$ factors, and
$\rk\mathcal B_k\le(D+2)^{k_{\mathrm{sl}}}$) has at most $|C_k|/(16(D+2)^{k_{\mathrm{sl}}})$ nonzero entries in every column $c$ that is
$k$-few.

Each $c\in Z_{\mathrm{bad}}$ is $k$-few for at least $(6K-1)-(3K-1)=3K$ indices $k$. So some $k$ has
$Z_k:=\{c\in Z_{\mathrm{bad}}:c\text{ is $k$-few}\}$ with $|Z_k|\ge3K|Z_{\mathrm{bad}}|/(6K-1)\ge|Z_{\mathrm{bad}}|/2$. Take rows $C_k$ and
columns $Z_k\subseteq C_\pi$ (disjoint). Lemma~\ref{lem:GL} gives a zero rectangle $C_k''\times Z_k''$ of $\mathcal B_k$ with halves.
There $[\dd xc^2]$ is covered by at most $k_{\mathrm{sl}}\le\kst$ types: column functions (the entries of $\Lambda_0(c)$, and
$\dd cs^2$), constants (the elements of $\Dist(S\cup B)$, $\beta(\pi)$ and $\beta(k)$), and row functions $\dd xs^2$.
Theorem~\ref{thm:LR} and Lemma~\ref{lem:typeOR} give a single-type rectangle $R\times C$ with $\min(|R|,|C|)\le h$ and
$|R|\ge2^{-\Lst}|C_k|/2>h$ (as $|C_k|\ge M\ge\sigma$). So $|C|\le h$, hence $|Z_k''|\le2^{\Lst}h$, $|Z_k|\le e$ and
$|Z_{\mathrm{bad}}|\le2e$.

\emph{(b) Choice of $c$.} Since $|C_\pi|\ge M\ge M_{\mathrm{sl}}>2\sigma\ge2e$, there is $c\in C_\pi\setminus Z_{\mathrm{bad}}$. Fix such a $c$
and a matching $\{(k,\lambda_k):k\in I_1\}$ of $G_c$ with $|I_1|=3K$; the $\lambda_k$ are distinct.

\emph{(c) Label conflicts.} Consider the digraph on $I_1$ with an arc $k\to k'$ whenever $k'\ne k$ and $\lambda_k\in\beta(k')$.
By (F3) each $\lambda_k$ equals at most one of the pairwise distinct numbers $\beta_i(k')$, so every out-degree is at most
$1$. The underlying graph therefore has at most $3K$ edges, hence average degree at most $2$, and
Theorem~\ref{thm:CW} gives an independent set of size at least $3K/3=K$. Let $I'$ consist of $K$ of its elements. For
distinct $k,k'\in I'$ we then have $\lambda_k\notin\beta(k')$.

\emph{(d) The new family.} Put $B':=B\cup\{c\}$, $\beta'(k):=(\beta(k),\lambda_k)$ and
\[
C'_k:=\bigl(C_k\cap\Sph(c,\lambda_k)\bigr)\setminus\bigcup_{s\in S}\Sph\bigl(s,\{\dd c{s'}:s'\in S\}\cup\{\lambda_{k'}:k'\in I'\}\bigr)\qquad(k\in I').
\]
Then $|C'_k|\ge\theta_{\mathrm{sl}}|C_k|-r(r+K)\sigma\ge\theta_{\mathrm{sl}}M-n(n+K)\sigma\ge\theta_{\mathrm{sl}}M/2$, because
$\theta_{\mathrm{sl}}M\ge\theta_{\mathrm{sl}}M_{\mathrm{sl}}=2n(n+K)\sigma$. We verify (F1)--(F5) for $\mathcal F'$.
\begin{itemize}[leftmargin=1.5em]
\item[(F1)] $c\notin S\cup B$, and $S\cup B\cup\{c\}$ is rainbow, by (F4) for $c\in C_\pi$.
\item[(F2)] $\dd x{b_i}=\beta_i(k)$ and $\dd xc=\lambda_k$ on $C'_k$.
\item[(F3)] The old values are pairwise distinct (F3), and the $\lambda_k$ are pairwise distinct (matching). For $k\in I'$,
  $\lambda_k\notin\beta(k)$ (edge rule) and $\lambda_k\notin\beta(k')$ for $k'\in I'\setminus\{k\}$ (independence). None of the
  values lies in $\Dist(S\cup B\cup\{c\})=\Phi(c)$: $\lambda_k\notin\Phi(c)$ by the edge rule; $\beta_i(k)\notin\Dist(S\cup B)$ by
  (F3); $\beta_i(k)\ne\beta_{i'}(\pi)$ by (F3), as $k\ne\pi$; and $\beta_i(k)\ne\dd cs$ by (F5) for $c$.
\item[(F4)] Let $x\in C'_k$. Then $x\notin S\cup B\cup\{c\}$, because $x\in C_k$, which is disjoint from $C_\pi\ni c$. The distances of
  $S\cup B\cup\{c,x\}$ are those of $S\cup B\cup\{x\}$ (rainbow), those of $S\cup B\cup\{c\}$ (rainbow), and $\dd xc=\lambda_k$.
  The cross coincidences are excluded as follows. $\dd xs=\dd c{s'}$ is excluded by the removal in the definition of
  $C'_k$. $\dd xs\in\beta(\pi)$ is excluded by (F5). $\beta_i(k)=\dd cs$ is excluded by (F5) for $c$. $\beta_i(k)=\beta_{i'}(\pi)$
  is excluded by (F3). $\lambda_k\in\Phi(c)$ and $\lambda_k\in\beta(k)$ are excluded by the edge rule. Finally $\lambda_k=\dd xs$
  is excluded by the removal.
\item[(F5)] For $x\in C'_k$ and $s\in S$: $\dd xs\notin\beta(I')$ by (F5) for $\mathcal F$, and $\dd xs\ne\lambda_{k'}$ by the
  removal.
\end{itemize}
So $\mathcal F'$ is a rainbow family with $p+1$ centres, $|I'|=K$ and $M(\mathcal F')\ge\theta_{\mathrm{sl}}M/2$.
\end{proof}

\subsection{The class-count theorem}

\begin{theorem}[E1$_r$]\label{thm:E1}
Let $0\le r\le n-5$, and write $K^{(p)}:=K^{(r,p)}$ for $1\le p\le n-4-r$ (Definition~\ref{def:typecounts}), so that
$K^{(1)}=K_r$ and $K^{(n-4-r)}=12$. For $1\le p\le n-5-r$ let $k_{\mathrm{sl}}(p):=k_{\mathrm{sl}}(r,p)$,
$\theta_{\mathrm{sl}}(p):=1/(16\Gamma'(D+2)^{k_{\mathrm{sl}}(p)})$ and $\hat M_{\mathrm{sl}}(p):=2n(n+K^{(p+1)})/\theta_{\mathrm{sl}}(p)$. Define
$\hat m_{n-4-r}:=\hat M_{\mathrm{core}}$ (Section~\ref{sec:core}) and, for $p=n-5-r,\dots,1$,
\[
\hat m_p:=\max\bigl(\hat M_{\mathrm{sl}}(p),\,2\hat m_{p+1}/\theta_{\mathrm{sl}}(p)\bigr),\qquad\text{and}\qquad
{\mathrm{ratio}}_r:=2\hat m_1+2n(n+K_r).
\]
These numbers depend only on $n$, $D$ and $\Gamma'$. If HR$_r$ holds with constant $\Gamma'$, then
$|\Bg_b({\mathrm{ratio}}_r\,\sigma)|<K_r$ for every $b\in W$.
\end{theorem}

\begin{proof}
Put $\tau:={\mathrm{ratio}}_r\,\sigma$, and suppose some $b\in W$ has $K_r$ radii $\rho_k$ ($k\in I$, $|I|=K_r$) in $\Bg_b(\tau)$.
Define $B:=\{b\}$, $\beta(k):=(\rho_k)$ and
\[
C_k:=\bigl(W\cap\Sph(b,\rho_k)\bigr)\setminus\bigcup_{s\in S}\Sph\bigl(s,\{\dd b{s'}:s'\in S\}\cup\{\rho_{k'}:k'\in I\}\bigr).
\]
Then $|C_k|\ge\tau-r(r+K_r)\sigma\ge\tau/2\ge\hat m_1\sigma$, because $\tau\ge2n(n+K_r)\sigma$. We check that this is a rainbow family.
(F1): $b\in W\subseteq X\setminus S$ and $S\cup\{b\}$ is rainbow. (F2) is clear. (F3): the $\rho_k$ are distinct and do not lie in
$R_b=\Dist(S\cup\{b\})$, by the definition of $\Bg_b$. (F4): for $x\in C_k$, $x\notin S\cup\{b\}$ and $S\cup\{x\}$ is rainbow (as
$x\in W$). Moreover $\dd xb=\rho_k\notin R_b$; $\dd xs\ne\dd b{s'}$ and $\dd xs\ne\rho_k=\dd xb$ by the removal; so
$S\cup\{b,x\}$ is rainbow by Lemma~\ref{lem:ext}. (F5) holds by the removal.

If $r=n-5$ then $p=1=n-4-r$, $K_r=12$ and $\hat m_1=\hat M_{\mathrm{core}}$, which contradicts Theorem~\ref{thm:core}. Otherwise, for
$p=1,\dots,n-5-r$ in turn, we have a family with $p$ centres, $K^{(p)}=6K^{(p+1)}$ classes and $M\ge\hat m_p\sigma\ge\hat M_{\mathrm{sl}}(p)\sigma$.
Lemma~\ref{lem:slicing}, with $K:=K^{(p+1)}$ and $k_{\mathrm{sl}}=k_{\mathrm{sl}}(r,p)\le\kst$, gives a family with $p+1$ centres,
$K^{(p+1)}$ classes and $M\ge\theta_{\mathrm{sl}}(p)\hat m_p\sigma/2\ge\hat m_{p+1}\sigma$. At $p=n-4-r$ we obtain $12$ classes and
$M\ge\hat M_{\mathrm{core}}\sigma=M_{\mathrm{core}}$, which contradicts Theorem~\ref{thm:core}.
\end{proof}

\begin{remark}[exceptional choices in a slicing round]
For the reader's orientation, these are the five ways in which a slicing round can go wrong, and how each is
handled. (i) $c\in Z_{\mathrm{bad}}$: at most $2e$ points, which are avoided. (ii) Small classes of points of $S$ through $c$
or with radius $\lambda_{k'}$: removed, at a cost of at most $r(r+K)\sigma$ points per class. (iii) Label conflicts
$\lambda_k\in\beta(k')$: handled by the independent set, at a loss of a factor $3$. (iv) Radii in $\Phi(c)\cup\beta(k)$:
these are never edges, and they enter the cover as constants or column functions. (v) Bisectors $\Bis(c,s)$: these
are row types. Only the covering form of the cleaning lemma is used, never a density-to-structure statement.
\end{remark}

% ====================================================================
\section{The threshold tower and the proof of the Main Theorem}\label{sec:main}
% ====================================================================

\subsection{The tower}
Let $n\ge4$ and $D\ge1$. The constants are defined in the following order.
\begin{enumerate}[label=(\arabic*),leftmargin=2em]
\item The type counts and class counts of Definition~\ref{def:typecounts}; these depend only on $n$.
\item The absolute parameters $\Lst,\mu,e,e_A,\smin$ of Definition~\ref{def:absolute}; these depend only on $n$, $D$ and
  $h=s_n(\lfloor D/2\rfloor)$.
\item $Q$ and $\Gamma_{n-2},\Gamma_{n-3},\Gamma_{n-4}$, by Theorem~\ref{thm:Rplus} and~\eqref{eq:Gamma-top}; these depend only on $n$
  and $D$.
\item For $n\ge5$ and $r=n-5,n-6,\dots,0$, in this order:
  \begin{equation}\label{eq:Gamma-low}
  \Gamma'_r:=\Gamma_{r+1}+r^2,\qquad{\mathrm{ratio}}_r:=\text{the number of Theorem~\ref{thm:E1} with }\Gamma'=\Gamma'_r,\qquad
  \Gamma_r:=P(K_r-1,r)\bigl(2\,{\mathrm{ratio}}_r\,\Gamma_{r+1}+2r^2+1\bigr),
  \end{equation}
  with $P(K,r)$ as in Lemma~\ref{lem:kill}. Here ${\mathrm{ratio}}_r$ depends only on $n$, $D$ and $\Gamma'_r$, and $\Gamma_r$ only on $n$,
  $D$ and $\Gamma_{r+1}$.
\end{enumerate}
No step refers to a later one. In particular $\smin$ involves no $\Gamma$, and no $\Gamma$ involves $\smin$ or $h$
(referee item~X5). The quantities $Q$ and $L_{k'}(D+2)\le\Lst$ enter only through $\Gamma_{n-3}$, $\Gamma_{n-4}$ and $\smin$
respectively.

\begin{theorem}\label{thm:allclaims}
For every $0\le j\le n-2$, Claim$(j)$ holds with the constant $\Gamma_j$ defined above.
\end{theorem}
\begin{proof}
For $j\ge n-4$ this is Theorem~\ref{thm:Rplus}. Let $n\ge5$ and $0\le r\le n-5$, and assume Claim$(r+1)$ with constant
$\Gamma_{r+1}$. Let $(S,W,\sigma)$ be a level-$r$ configuration with $\sigma\ge\smin$. By Lemma~\ref{lem:HR}, HR$_r$ holds for
$(S,W,\sigma)$ with constant $\Gamma'_r$. By Theorem~\ref{thm:E1}, $|\Bg_b({\mathrm{ratio}}_r\sigma)|\le K_r-1$ for every
$b\in W$. Lemma~\ref{lem:kill}, with $K=K_r-1$ and $\tau={\mathrm{ratio}}_r\sigma\ge\sigma$, together with $1\le\sigma$
(Remark~\ref{rem:killconst}), gives $|W|\le P(K_r-1,r)(2{\mathrm{ratio}}_r\Gamma_{r+1}+2r^2+1)\sigma=\Gamma_r\sigma$. Downward
induction on $r$ completes the proof.
\end{proof}

The induction is legitimate because each Claim$(j)$ is a statement about all configurations and all caps
$\sigma\ge\smin$, with a constant independent of $\sigma$. The proofs apply Claim$(j+1)$ with caps
$\theta\ge\sigma$ (Lemma~\ref{lem:HR}, and through it Lemma~\ref{lem:step3} and Lemma~\ref{lem:slicing}(a)), $Q\sigma$
(Theorem~\ref{thm:Rplus}) and ${\mathrm{ratio}}_r\sigma$ (Lemma~\ref{lem:kill}). Each of these caps is at least $\smin$.

\subsection{Size of the constants}

\begin{lemma}[growth]\label{lem:growth}
Put $a:=D+2$. There is $c_n$, depending only on $n$, with $\Gamma_0\le a^{c_n}$ for all $D\ge1$. More precisely,
$g_r:=\log_a\Gamma_r$ satisfies $g_r\le(n-2-r)\,g_{r+1}+b_n$ for $0\le r\le n-5$, where $b_n$ depends only on $n$. Hence
$g_0\le\frac{(n-2)!}2\,(g_{n-4}+b_n)$, and $g_{n-4}=O(n^2)$.
\end{lemma}
\begin{proof}
Call a quantity \emph{$n$-polynomial} if it is at most $a^c$ with $c$ depending only on $n$. Since $a\ge3$, every number
depending only on $n$ is $n$-polynomial, and sums, products and maxima of boundedly many (in terms of $n$)
$n$-polynomial quantities are $n$-polynomial. $Q$ is $n$-polynomial, and so are $P^+_j\le a^{2e_j}$ and
$P(K,r)\le a^{2e_{K,r}}$ (Lemma~\ref{lem:DIM}). Hence $\Gamma_{n-2},\Gamma_{n-3},\Gamma_{n-4}$ are $n$-polynomial. Since $e_j=O(n^2)$, $k'=O(n^2)$ and
$\log_a(2^{44}n^4)=O(\log n)$, two applications of~\eqref{eq:Gamma-top} give $g_{n-4}=O(n^2)$.

For the recursion, fix $r\le n-5$ and write $\Gamma'=\Gamma'_r$. First, $\Gamma_{r+1}\ge r^2$: this is clear from~\eqref{eq:Gamma-top}
and~\eqref{eq:Gamma-low}. Hence $\Gamma'\le2\Gamma_{r+1}$. Next, $2/\theta_{\mathrm{sl}}(p)=32\Gamma'a^{k_{\mathrm{sl}}(p)}$ and
$\hat M_{\mathrm{sl}}(p)=32\Gamma'a^{k_{\mathrm{sl}}(p)}n(n+K^{(p+1)})$. Using $\max(x,y)\le xy$ for $x,y\ge1$, we get
$\hat m_p\le32\Gamma'a^{k_{\mathrm{sl}}(p)}n(n+K^{(p+1)})\,\hat m_{p+1}$, and hence
\[
\hat m_1\le\Bigl[\prod_{p=1}^{n-5-r}32\Gamma'a^{k_{\mathrm{sl}}(p)}n(n+K^{(p+1)})\Bigr]\,2^{20}4^nn^2\Gamma'^2a^{2k_c(r)}
\le A_n\,a^{B_n}\,\Gamma'^{\,n-3-r},
\]
with $A_n,B_n$ depending only on $n$. Also $\hat m_1\ge2n(n+K_r)$, so ${\mathrm{ratio}}_r\le3\hat m_1$. Finally
$2r^2+1\le{\mathrm{ratio}}_r\Gamma_{r+1}$, so~\eqref{eq:Gamma-low} gives
$\Gamma_r\le3P(K_r-1,r)\,{\mathrm{ratio}}_r\,\Gamma_{r+1}\le9a^{2e_{K_r-1,r}}A_na^{B_n}2^{n}\Gamma_{r+1}^{\,n-2-r}$. Taking $\log_a$
proves the recursion, and iterating it from $r=n-5$ down to $r=0$ gives the bound on $g_0$, since
$\prod_{r=0}^{n-5}(n-2-r)=(n-2)!/2$.
\end{proof}

By Definition~\ref{def:typecounts}, $\kst\le6^{n-4}+2n^2$. So $\Lst=L_{\kst}(D+2)=O(6^{3n}\sqrt D\log D)$ and
\[
\smin=2^{\Lst+2+2n}(h+1)=2^{O(6^{3n}\sqrt D\log D)}\,(h+1).
\]
We have made no attempt to optimise the constants: the exponent $c_n$ grows like $(n-2)!$ and the constant in
$\Lst$ like $6^{3n}$. Only the shape $2^{O_n(\sqrt d\log d)}$ is claimed.

\subsection{Proof of the Main Theorem}

\begin{theorem}[Main Theorem]\label{thm:main}
For every $n\ge3$ there is $C_n$ such that $f_d(n)\le2^{C_n\sqrt d\log_2(d+2)}$ for all $d\ge1$.
\end{theorem}
\begin{proof}
For $n=3$ this follows from Blokhuis's theorem: rainbow-$3$-free sets are exactly the isosceles sets, including
degenerate triples, and these have at most $\binom{d+2}2$ points~\cite{Blokhuis}. Hence $f_d(3)\le\binom{d+2}2+1$.

Let $n\ge4$ and $D\ge1$, and let $X$ be a finite rainbow-$n$-free set of affine dimension at most $D$. We may assume
$X\subset\R^D$. The triple $(\emptyset,X,\smin)$ is a level-$0$ configuration with $\smin\ge\smin$, so
Theorem~\ref{thm:allclaims} (Claim$(0)$) gives
\[
|X|\le\Gamma_0\,\smin=\Gamma_0\,2^{\Lst+2+2n}\,(h+1).
\]
By induction on $D$, $h=s_n(\lfloor D/2\rfloor)$ is finite (with $s_n(0)=1$). So $s_n(D)$ is finite and
$s_n(D)\le A(D)\,(s_n(\lfloor D/2\rfloor)+1)$ with $A(D):=\Gamma_0\,2^{\Lst+2+2n}$. By Lemma~\ref{lem:growth} and the definition
of $L_k$, $\log_2A(D)\le c_n\log_2(D+2)+2^{14}\kst^{\,3}\sqrt{D+4}\log_2(D+4)+2+2n$. For $D\ge1$ we have
$\sqrt{D+4}\le\sqrt5\sqrt D$, $\log_2(D+4)\le2\log_2(D+2)$ and $\sqrt D\log_2(D+2)\ge1$, so
$A(D)\le2^{C\sqrt D\log_2(D+2)}$ for a constant $C\ge1$ depending only on $n$. Lemma~\ref{lem:halving} gives
$s_n(d)\le2^{6C\sqrt d\log_2(d+2)}$, and $f_d(n)=s_n(d)+1\le2^{(6C+1)\sqrt d\log_2(d+2)}$.
\end{proof}

\begin{remark}
The proof uses Theorem~\ref{thm:LR} in five places: in G1 (through $\mu$); in Step~A of Theorem~\ref{thm:Rplus}; in
Steps~3--5 of the Core Lemma; and in Lemma~\ref{lem:slicing}(a). If Theorem~\ref{thm:LR} were replaced by the trivial
bound $2^{-O(r)}$ for single-type rectangles, the same argument would give only $s_n(D)\le2^{O(D)}s_n(\lfloor D/2\rfloor)$, that
is, $f_d(n)=2^{O_n(d)}$, the classical order of magnitude.
\end{remark}

% ====================================================================
\section{Concluding remarks}\label{sec:concl}
% ====================================================================

\subsection{What the argument rests on}
Apart from classical results (John, Sheppard, Caro--Wei, K\"onig, Blokhuis), every step is proved in this manuscript. The
load-bearing new ingredient is Theorem~\ref{thm:LR}. The geometric induction (Sections~\ref{sec:framework}--\ref{sec:main}) uses
it only as a black box, and all other inputs are elementary. In decreasing order of what we regard as the risk of an error,
the parts that most deserve an expert's attention are the following.
\begin{enumerate}[leftmargin=2em]
\item \emph{Theorem~LR}, in particular the normal form (Lemma~\ref{lem:NF}) and the rescaling at the largest gap by a row
  function, a column function or a constant. Its constants-only and column-only special cases are close to known
  arguments~\cite{Lovett,Rothvoss}. The mixed row/column case rests on Lemmas~\ref{lem:A}, \ref{lem:B} and~\ref{lem:NF}.
\item \emph{Step~A} of Theorem~\ref{thm:Rplus} (Lemma~\ref{lem:stepA}). This is the only place where a density statement (the
  cross-pair count~\eqref{eq:S2}) is converted into a covering statement, at the cost of a fourth root
  (Lemma~\ref{lem:twodist}). An earlier version contained a genuine threshold gap here (fix~F1).
\item The bookkeeping of the Core and Slicing Lemmas, in particular the completeness of the coincidence lists in
  Lemmas~\ref{lem:step0}, \ref{lem:step1} and~\ref{lem:51}, and the verification of (F1)--(F5) in Lemma~\ref{lem:slicing}(d).
\end{enumerate}
If Theorem~\ref{thm:LR} failed, the argument would give only the classical bound $2^{O(d)}$.

\subsection{Computational sanity checks}
Several finite, threshold-free parts of the argument were tested by computer, in the project and by the referees.
We list these tests for completeness. They are \emph{sanity checks only}. Every threshold in the argument
($\smin\ge2^{\Lst}(h+1)$, the cleaning densities, the sides required in the orthogonal splits) is astronomically
larger than any set that can be handled numerically, so no threshold-dependent conclusion was, or could be,
tested. Only inclusions and implications were tested.
\begin{itemize}[leftmargin=1.5em]
\item \emph{Rounding and cleaning} (referee pass on Theorem~LR). John positions were computed numerically (as
  minimum-volume enclosing ellipsoids) on $k$-listable instances, giving $\max|u_x|/s^{1/4}\le1.0002$ (the solver
  tolerance) and $\max|v_w|/(Ks^{1/4})\le1$. Sheppard's formula was compared with Monte Carlo estimates and agreed to
  within $10^{-3}$. The cleaning lemma was tested on adversarial sparse low-rank matrices, and a zero block was always
  found.
\item \emph{Normal form} (same referee pass). An independent implementation of Lemmas~A, B, NF and one round of the LR
  induction was run on $1600$ adversarial block-structured instances: accumulating labels $\kappa\pm2^{-j}$ with $j$ up
  to $40$, ties between row and column functions on blocks, magnitudes $2^{\pm40}$, and up to $5$ types of mixed kinds.
  None of (NF1)--(NF3), the bounds on $|N|$, the identification of the minority set with the support of $H$, or the
  rank bound for $H$ was ever violated.
\item \emph{Genericity in Step~A} (referee pass on Theorem~R$^+$). In $70{,}213$ random grid configurations
  ($n=4,5,6$), $20{,}361$ of them non-generic, no non-generic point fell outside the listed loci.
\item \emph{The case $n=5$} (referee pass on the $n=5$ argument). About $150{,}000$ exact-arithmetic checks of the
  threshold-free inclusions of Steps~0--3 and~5 were run on rainbow-$5$-free families (Johnson-type schemes, products, and
  greedy grid sets), with no violation.
\item \emph{General $n$, author's checks.} $348{,}884$ exact-integer checks for $n=6,7$ and all relevant $(r,p)$ were run,
  covering the non-generic loci of Step~0, the inclusions of Step~1, the characterisation in Lemma~\ref{lem:51}, the
  equivalence behind Lemma~\ref{lem:HR}, and the family conditions in Lemma~\ref{lem:slicing}(d). There was no violation.
  These checks sampled \emph{arbitrary} integer points, not points of rainbow-$n$-free sets (referee item~X6). So they
  do not test the one implication that uses rainbow-$n$-freeness, namely the label lemma behind G1.
\item \emph{General $n$, independent checks} (referee pass on the general-$n$ argument). About $340{,}000$ exact checks
  were run on genuinely rainbow-$n$-free sets for $n=6,7$. The sets were orthogonal unions of integer points on small
  spheres, arranged so that cross distances collide, together with few-distance lattices. The checks included the
  full pipeline of Lemmas~\ref{lem:HR} and~\ref{lem:kill}, the initial family of Theorem~\ref{thm:E1}, slicing rounds with
  actual maximum matchings and K\"onig covers, and Steps~0--2 and~5 of the Core Lemma. They also included the label
  lemma (Lemma~\ref{lem:label}) on label-disjoint generic pairs. There was no violation. On these sets families die
  after one or two slicing rounds, and at most $7$ (for $n=6$) or $12$ (for $n=7$) classes were available. So the
  contradiction in the Core Lemma itself was never exercised; only its steps were.
\end{itemize}

\subsection{Literature}
In the referee pass on Theorem~LR, the following works were searched for prior statements of
Theorem~\ref{thm:LR} or Corollary~\ref{cor:RQlist}:~\cite{HMST}, which proves bounds for nonnegative
integer matrices with bounded average and remarks that binarity is not crucial; Sudakov--Tomon~\cite{SudakovTomon};
Rothvoss~\cite{Rothvoss}; Lovett~\cite{Lovett}; and works citing~\cite{SingerSudan}. No statement of a
$2^{-\tilde O(\sqrt d)}$ bound for $k$-listable real matrices was found. The search was not exhaustive. As stressed in
Remark~\ref{rem:SS}, the possible improvement over the bound stated in~\cite{SingerSudan} for parallel partitions should
be confirmed by experts. Corollary~\ref{cor:RQlist} follows from Rothvoss's argument in a few lines, and it may be regarded
as folklore.

\subsection{Further remarks}
\begin{itemize}[leftmargin=1.5em]
\item The class counts $K_r=12\cdot6^{n-5-r}$, the type counts and all constants are far from optimal. For $n=5$ a referee
  observed that the count in Step~5 already works with $7$ classes instead of $12$, at the cost of larger constants.
\item The super-polynomial factor $2^{\Lst}$ in the thresholds cannot be removed from the class-count statement E1$_r$
  alone. If $Y$ is rainbow-$(n-1)$-free and $o$ is any point, then $Y\cup\{o\}$ is rainbow-$n$-free. A referee observed
  that, for $n=5$, taking $Y$ to be a Johnson-type $5$-distance set and $o$ a suitable weight vector gives a point
  with many more than $12$ classes, each of size polynomial in $D$. So E1 with a threshold polynomial in $D$ would be
  false. With the thresholds used here, E1 is consistent with every example known to us.
\item No super-polynomial lower bound for $f_d(n)$ is known to us, so the gap between $d^{\Omega_n(1)}$ and
  $2^{O_n(\sqrt d\log d)}$ remains wide.
\end{itemize}

% ====================================================================
\section{Formal verification in Lean}\label{sec:lean}
% ====================================================================
\begin{sloppypar}

The complete argument of this paper has been formalized in Lean~4 (toolchain \texttt{v4.33.1}) on Mathlib revision
\texttt{0df444a3}, the same pins as the Google DeepMind \emph{formal-conjectures} repository at commit
\texttt{c77fdd446d23539302aa73bcfae4a9fa962b8382}. The formalization is about $9{,}800$ lines.

\paragraph{The target statement.} The theorem proved is \texttt{Erdos1088.erdos\_1088}:
\[
\forall n\ge 3,\qquad \bigl(d\mapsto \log_2 f(d,n)\bigr) = o(d)\quad (d\to\infty),
\]
where \texttt{f d n} is defined as $\inf\{m\ge n:$ every $m$-point subset of $\R^d$ contains an $n$-point subset
with $\binom n2$ distinct distances$\}$. The statement and all definitions it uses (\texttt{f}, \texttt{cardSet},
\texttt{HasSubsetWithDistinctDistances}, \texttt{PairwiseDistancesDistinct}, \texttt{distinctDistances},
\texttt{distanceSet} and the notation $\R^d$) are copied character for character from formal-conjectures. A script
checks this against the upstream files. Only module headers and metadata attributes are dropped.

Lean's \texttt{sInf} of an empty set of naturals is $0$, so a proof could in principle exploit an empty
\texttt{cardSet}. The formal proof does not do this: it exhibits an explicit element of \texttt{cardSet d n} and
proves $f(d,n)\le n\cdot 2^{C_n\sqrt d\log_2(d+2)}$, which is the Main Theorem.

\paragraph{Trust base.} \texttt{\#print axioms Erdos1088.erdos\_1088} reports only \texttt{propext},
\texttt{Classical.choice} and \texttt{Quot.sound}, the standard axioms of Lean and Mathlib. The development
contains no \texttt{sorry}, no \texttt{axiom} declarations and no \texttt{native\_decide}.

\paragraph{Differences from the informal proof.} The formalization follows Sections~\ref{sec:LR}--\ref{sec:main}
with the following changes, none of which affects the Main Theorem.
\begin{itemize}[leftmargin=1.5em]
\item The normalization step (NF3) of Lemma~\ref{lem:NF} is replaced by a tournament argument. As a result, the formal
  version of Theorem~\ref{thm:LR} has loss $2^{O(k)}\sqrt r\log r$ instead of $O(k^3\sqrt r\log r)$; for fixed $k$ the
  shape $2^{-O_k(\sqrt r\log r)}$ is the same, and this is all the later sections use.
\item John's theorem in Lemma~\ref{lem:R} is replaced by a factorization derived from the Kiefer--Wolfowitz
  optimal-design theorem, proved from the matrix determinant lemma and the spectral theorem. Sheppard's formula is
  proved from the rotation invariance of the standard Gaussian.
\item Several constants are larger. These include those in the Core Lemma and in Theorem~\ref{thm:Rplus}, the
  number of big radii at which Step~A applies ($3$ instead of $2$), and $\kst$. The case $n=3$ is derived from G1
  rather than from Blokhuis's theorem.
\item Lean made explicit some hypotheses that the text leaves implicit, for example $\sigma>0$ and $\Gamma\ge 0$ in
  the kill step, and $B\subseteq X$ in the Core Lemma. It found no gap in the argument.
\end{itemize}

\paragraph{Availability.} The Lean source, build instructions and audit scripts are available at
\url{https://github.com/H0W4RD-dev/erdos-1088-lean}.
\end{sloppypar}

% ====================================================================
\section{AI assistance and verification; credits}\label{sec:ai}
% ====================================================================

\subsection*{AI assistance and verification}
The mathematics in this manuscript was developed by AI agents (Claude, Anthropic) in research sessions directed by the
team. The agents proposed the strategy, proved the lemmas, and wrote the project notes from which this manuscript was
assembled. The manuscript itself was also prepared by an AI agent.

The arguments were checked by \emph{independent AI referee passes}: separate agent sessions that had not written the
material under review, re-derived each step by hand, and ran their own computer tests. In summary:
\begin{itemize}[leftmargin=1.5em]
\item one pass on Theorem~\ref{thm:LR} (with Lemma~\ref{lem:R}, the cleaning lemma, Lemmas~A, B and NF), on
  Corollary~\ref{cor:RQlist} and its relation to~\cite{SingerSudan}, and on an earlier proof of the case $n=4$. Verdict:
  correct, with editorial remarks only;
\item one pass on the kill step, the two-class count and Theorem~R$^+$. Verdict: correct after fixes F1--F3, of which
  F1 repaired a genuine gap;
\item one pass on the case $n=5$ of the class-count argument. Verdict: correct with expository gaps; items A--G of that
  report are addressed in Sections~\ref{sec:framework}--\ref{sec:core};
\item one pass on the general-$n$ argument (configurations, HR, kill step, Core Lemma, Slicing Lemma, threshold tower).
  Verdict: correct with expository gaps; items X1--X6 of that report are addressed in Sections~\ref{sec:prelim},
  \ref{sec:framework}--\ref{sec:main} and~\ref{sec:concl}.
\end{itemize}
Each referee found the material under review correct, conditional on its inputs. All of them identified
Theorem~\ref{thm:LR}, which is new and unpublished, as the weakest link.

\textbf{No human expert has reviewed any part of this argument.} AI referee passes can miss errors, and several
passes sharing similar blind spots is a real possibility. The Lean formalization of Section~\ref{sec:lean} was also
written by AI agents. It removes the dependence on the referees for the \emph{formal} statement: Lean's kernel checks
every step. Two things remain for human judgement: whether the formal statement faithfully captures Erd\H{o}s's
question, and the sharper form of Theorem~\ref{thm:LR} stated here, whose NF3 step is not formalized.

We also note what has \emph{not} been refereed in the present form. While consolidating the project notes into this
manuscript we made some routine changes, and these have not been checked by any referee pass. They are: the
explicit choices $Q=2^{44}n^4(D+2)^{8k'}$, $\smin=2^{\Lst+2+2n}(h+1)$ and $e_A=2^{\Lst+2}(h+1)$; the explicit constant
$2^{14}$ in $L_k(r)$ and the explicit cost bounds in Lemma~\ref{lem:NF}; the column-basis proof of the cleaning lemma; the
self-contained statement and proof of Lemma~\ref{lem:twodist}; the direct (rather than translated) single-cap proof of
Theorem~\ref{thm:Rplus}; and the explicit growth recursion in Lemma~\ref{lem:growth}. We believe these changes are correct
and harmless, but they are part of what an expert should check.

\subsection*{Author list and credits}
\textbf{[PLACEHOLDER --- the author list and the form of credit are to be confirmed by the team.]}

\smallskip\noindent
\emph{Direction.} Will Howe (H0W4RD-dev) directed the work.

\smallskip\noindent
\emph{Contributors and acknowledgements.} The team also includes Jensen Kohlmeyer (jenw1n), Liam Kruer (lkruer) and Grant
Huffman (Ghuffy1). They are listed here as team members; the project records available to us do not attribute any
mathematical contribution to this result to them. We acknowledge the prior in-project work of a peer session~\cite{Peer},
on which Section~\ref{sec:intro} and Lemma~\ref{lem:H} build.

% ====================================================================
\begin{thebibliography}{99}
% ====================================================================

\bibitem{Ball}
K.~Ball, An elementary introduction to modern convex geometry, in: \emph{Flavors of Geometry}, MSRI Publications~31,
Cambridge University Press, 1997, pp.~1--58.

\bibitem{Blokhuis}
A.~Blokhuis, \emph{Few-distance sets}, CWI Tract~7, Centrum voor Wiskunde en Informatica, Amsterdam, 1984.

\bibitem{Caro}
Y.~Caro, New results on the independence number, Technical Report, Tel Aviv University, 1979.

\bibitem{EP1088}
T.~F.~Bloom, Erd\H{o}s Problem \#1088, \url{https://www.erdosproblems.com/1088}. The problem page records the
Erd\H{o}s--Straus bound $f_d(n)\le c_n^d$ and its source.

\bibitem{GavinskyLovett}
D.~Gavinsky and S.~Lovett, En route to the log-rank conjecture: new reductions and equivalent formulations, in:
\emph{Automata, Languages, and Programming (ICALP 2014)}, Lecture Notes in Computer Science~8572, Springer, 2014,
pp.~514--524.

\bibitem{GW}
M.~X.~Goemans and D.~P.~Williamson, Improved approximation algorithms for maximum cut and satisfiability problems using
semidefinite programming, \emph{J.~ACM} \textbf{42} (1995), 1115--1145.

\bibitem{HMST}
Z.~Hunter, A.~Milojevi\'c, B.~Sudakov and I.~Tomon, Disjoint pairs in set systems and combinatorics of low rank matrices,
arXiv:2411.13510.

\bibitem{John}
F.~John, Extremum problems with inequalities as subsidiary conditions, in: \emph{Studies and Essays Presented to
R.~Courant on his 60th Birthday}, Interscience, New York, 1948, pp.~187--204.

\bibitem{Konig}
D.~K\H{o}nig, Gr\'afok \'es m\'atrixok, \emph{Matematikai \'es Fizikai Lapok} \textbf{38} (1931), 116--119.

\bibitem{Lovett}
S.~Lovett, Communication is bounded by root of rank, \emph{J.~ACM} \textbf{63} (2016), Article~1; arXiv:1306.1877.

\bibitem{Peer}
Prior in-project work: research notes on Erd\H{o}s Problem \#1088 by a peer AI session of the same project
(notes file \texttt{e1088.md}, 2026), unpublished.

\bibitem{Rothvoss}
T.~Rothvo\ss, A direct proof for Lovett's bound on the communication complexity of low rank matrices,
arXiv:1409.6366, 2014.

\bibitem{Sheppard}
W.~F.~Sheppard, On the application of the theory of error to cases of normal distribution and normal correlation,
\emph{Philos. Trans. Roy. Soc. London Ser.~A} \textbf{192} (1899), 101--167.

\bibitem{SingerSudan}
N.~Singer and M.~Sudan, Point-hyperplane incidence geometry and the log-rank conjecture, arXiv:2101.09592 (version~5,
2022).

\bibitem{SudakovTomon}
B.~Sudakov and I.~Tomon, Matrix discrepancy and the log-rank conjecture, arXiv:2311.18524.

\bibitem{Wei}
V.~K.~Wei, A lower bound on the stability number of a simple graph, Bell Laboratories Technical Memorandum
81-11217-9, 1981.

\end{thebibliography}


\end{document}
